\documentclass[11pt,a4paper]{article}
\usepackage[T1]{fontenc}
\usepackage[utf8]{inputenc}
\usepackage{lmodern}
\usepackage[margin=26mm,headheight=14pt]{geometry}
\usepackage{amsmath,amssymb,amsthm,mathtools,mathrsfs}
\usepackage{microtype,booktabs,longtable,array}
\usepackage[dvipsnames]{xcolor}
\usepackage{enumitem}
\usepackage{xurl}
\usepackage{fancyhdr}
\usepackage{tocloft}
\setlength{\cftbeforesecskip}{4pt}
\setlength{\cftbeforetoctitleskip}{0pt}
\setlength{\cftaftertoctitleskip}{12pt}
\usepackage[colorlinks=true,linkcolor=MidnightBlue,citecolor=MidnightBlue,
 urlcolor=MidnightBlue,bookmarksnumbered=true]{hyperref}
\hypersetup{
 pdftitle={Coincident Stieltjes Products and Directional Zeta Identities},
 pdfauthor={Research continuation prepared for Vladimir Reshetnikov},
 pdfsubject={Finite parts, Tornheim derivatives, polylogarithms, harmonic sums and generalized Stieltjes constants},
 pdfkeywords={Hurwitz zeta, Tornheim zeta, Stieltjes constants, finite part, polylogarithm, Gamma, harmonic numbers}
}
\newtheorem{theorem}{Theorem}[section]
\newtheorem{proposition}[theorem]{Proposition}
\newtheorem{lemma}[theorem]{Lemma}
\newtheorem{corollary}[theorem]{Corollary}
\theoremstyle{definition}
\newtheorem{definition}[theorem]{Definition}
\theoremstyle{remark}
\newtheorem{remark}[theorem]{Remark}
\DeclareMathOperator*{\FP}{FP}
\DeclareMathOperator{\CT}{CT}
\DeclareMathOperator{\Li}{Li}
\numberwithin{equation}{section}
\allowdisplaybreaks[2]
\setlist{itemsep=3pt,topsep=5pt}
\setlength{\emergencystretch}{2em}
\pagestyle{fancy}
\fancyhf{}
\fancyhead[L]{\small Coincident Stieltjes products}
\fancyhead[R]{\small Directional zeta identities}
\fancyfoot[C]{\thepage}
\renewcommand{\headrulewidth}{0.3pt}

\begin{document}
\hypersetup{pageanchor=false}
\begin{titlepage}
\thispagestyle{empty}
\vspace*{16mm}
{\color{MidnightBlue}\Large ProveIt research continuation\par}
\vspace{12mm}
{\huge\bfseries Coincident Stieltjes Products\\[3mm]
and Directional Zeta Identities\par}
\vspace{8mm}
{\Large Cubic finite parts, distributional contact terms,\\[2mm]
coordinate invariance, and exact spectral reductions\par}
\vspace{13mm}
{\large Prepared for Vladimir Reshetnikov\par}
\vspace{3mm}
10 October 2026
\vspace{12mm}
\begin{abstract}
\noindent
This article continues the polylogarithm manuscript and six incoming
research archives at a pinned ProveIt revision. We prove a finite
holomorphic completion rule for coincident products of arbitrary
Stieltjes argument derivatives, give its explicit cubic Tornheim
realization, and reduce the requested cubic digamma finite part to the
third derivative of the diagonal Tornheim function. The same local
calculation evaluates its second derivative and, together with a
Mellin decomposition, yields exact quadratic derivatives on every
admissible asymmetric ray through the spectral origin.

We also determine all delta-supported corrections between periodic
distributional correlations and canonical finite parts in the shift
variable. A closed Gamma and harmonic-polynomial generator gives every
Stieltjes index and every argument-derivative order. We classify all
finite logarithmic singular germs invariant under coordinates with a
prescribed identity jet, including necessity and sharp Stieltjes
thresholds. For colored weighted sums, finite rational generating
functions give all outer-order derivatives in ordinary polylogarithms,
with exact Gamma and generalized-Stieltjes specializations. Finally,
we give all principal Laurent coefficients and finite parts at a
common-shift harmonic resonance with two independently moving
exponents and arbitrary trailing depth.

Proofs specify the cutoff coordinate, the meromorphic restriction, and
every allowed derivative direction. Classical ingredients and inherited
results are credited; no general historical-priority claim is made.
The Gaussian $S_6$ and current $S_8$ candidates remain conjectural.
Reproducible exact and high-precision checks, narrowly scoped source
corrections, and a proposed research agenda accompany the article.
\end{abstract}
\vfill
\noindent\small
Source revision:
\href{https://github.com/VladimirReshetnikov/ProveIt/tree/dcd95baeb7c0e914b138bddef6829c5b1d52b726}
{\texttt{dcd95baeb7c0e914b138bddef6829c5b1d52b726}}.
The accompanying source receipt records verified Git blob identities.
\end{titlepage}

\hypersetup{pageanchor=true}
\pagenumbering{roman}
{\small\tableofcontents}
\clearpage
\pagenumbering{arabic}
% BEGIN sections/01_scope.tex
\section{Scope, results, and conventions}
\label{scope:sec}

\subsection{The research questions being continued}

The source corpus is the manuscript in
\texttt{Analysis/Polylogarithms/docs/manuscript} and the six research
archives in \texttt{docs/incoming}, frozen at the revision on the title
page \cite{Repo}. A read-only extraction verified the Git blob identity
of 704 relevant text, code, data, bibliography, and archive files.
This is a source-integrity statement, not a claim to have verified every
theorem in those files. The mathematical audit is targeted to the
questions and premises used below.

Three developments in the incoming material make a further exact
calculus possible. First, periodic Stieltjes functions have been given
consistent distributional and pointwise finite-part conventions.
Second, nonlinear coordinate changes have been reduced to finite
residue expressions. Third, the shifted height-one harmonic family
has an exact Gamma-ratio representation with a finite Stirling
specialization. The present work uses those developments as
established inputs and extends the questions they explicitly leave
open \cite{IncomingCoincident,IncomingTransport,IncomingExact}.

The main outcomes and their precise scope are as follows.

\begin{longtable}{@{}>{\raggedright\arraybackslash}p{0.25\textwidth}
 >{\raggedright\arraybackslash}p{0.48\textwidth}
 >{\raggedright\arraybackslash}p{0.20\textwidth}@{}}
\toprule
Incoming question or target & Result established here & Location\\
\midrule
\endhead
Coincident products beyond degree two &
A finite local completion for every number of factors; a fully
explicit cubic global kernel at all argument orders &
Theorem \ref{cub:general-completion};
Proposition \ref{cub:global}\\[3pt]
Cubic digamma finite part &
An exact reduction to $\omega'''(0)$, where
$\omega(t)=T(t,t,t)$; also an exact evaluation of $\omega''(0)$ &
Theorem \ref{cub:diagonal}\\[3pt]
Contact terms at a shift collision &
One explicit Gamma--harmonic generator for every delta coefficient;
no undetermined lower delta derivatives &
Theorem \ref{cont:main}\\[3pt]
Coordinate-invariant singular germs &
Necessary and sufficient coefficient conditions for every prescribed
order of tangency &
Theorems \ref{invg:classification},
\ref{invg:stieltjes-threshold}\\[3pt]
Normal Tornheim derivatives at integer inner orders &
Finite depth-two reductions at all integer inner orders, and
depth-one closure at nonpositive orders with arbitrary weights &
Theorems \ref{sp:integer}, \ref{sp:rationaltheorem}\\[3pt]
Spectral direction dependence &
Constant, linear, and quadratic terms on every ray with
$(a+c)(b+c)\ne0$ &
Theorem \ref{sp:raytheorem}\\[3pt]
Independent harmonic exponents &
Complete Laurent principal parts and finite parts for two moving
exponents at every trailing depth and common shift &
Theorem \ref{h2:thm:direction}\\
\bottomrule
\end{longtable}

These statements answer several archived research questions, but they
do not all have the same kind of closure. The cubic finite part is
reduced to a specified multizeta derivative, not evaluated in ordinary
zeta constants alone. The positive-inner-order colored formulas close
in depth-two polylogarithms. The weighted nonpositive-inner-order
theorem closes in ordinary polylogarithms and their order derivatives.
The two-exponent harmonic theorem eliminates its holomorphic
remainder through the finite part, but does not eliminate the
remainder from every higher regular coefficient.

\subsection{Relation to the literature}

The analytic ingredients are classical. Hurwitz Fourier expansions
and product integrals are treated by Espinosa and Moll
\cite{EM2002,EM2006}. Bailey and Borwein systematically study
Mordell--Tornheim order derivatives and their polylogarithmic integral
representations \cite{BB2014}. Colored Tornheim partial fractions already appear
in Zhao \cite{sp:Zhao}. Borwein and Dilcher give direction-dependent
Tornheim values and first derivatives on the equal-inner-order family
\cite{BorweinDilcher}; related Witten values are discussed by Romik
\cite{sp:Romik}. Logarithmic finite parts and their change of variables
have an established theory \cite{Latt2011}. Meromorphic continuation
and Laurent algorithms for multiple zeta functions at integer points
are supplied by Akiyama--Egami--Tanigawa and
Matsumoto--Onozuka--Wakabayashi \cite{h2:AET,h2:MOW}.

Accordingly, this article makes no claim to have discovered finite
parts, the Tornheim--Hurwitz relation, directional regularization, or
multiple-zeta continuation. The contributions are the stated explicit
completion and contact generators, the invariant-subspace
classification, the asymmetric quadratic formula, and the finite
two-exponent harmonic coefficient rule relative to the audited corpus.
The elementary weighted identities are included as proved, useful
consequences of a general rational coefficient method.

\subsection{Notation and finite-part prescriptions}

We use
\begin{equation}\label{scope:stieltjes}
 \zeta(1+u,a)=\frac1u+
       \sum_{n\ge0}\frac{(-1)^n}{n!}\gamma_n(a)u^n,
 \qquad \gamma_0(a)=-\psi(a),
 \qquad \gamma_n=\gamma_n(1).
\end{equation}
Here $\psi=\Gamma'/\Gamma$, and $\gamma=\gamma_0$ is Euler's
constant. A prime on $\psi$ is an argument derivative. On $\zeta(s,a)$,
a prime or $[k]$ superscript denotes a derivative in $s$ unless an
argument derivative is explicitly written. Thus
$\zeta^{[k]}(s,a)=\partial_s^k\zeta(s,a)$ and
$\Li_s^{[k]}(z)=\partial_s^k\Li_s(z)$.
All integer summation bases are raised to complex powers using their
real logarithms. A common shift $a$ is initially in $\Re a>0$, with the
analytic logarithm in that half-plane.

The generalized harmonic numbers are
\[
 H_N^{(j)}=\sum_{\nu=1}^N\nu^{-j},\qquad H_N=H_N^{(1)},\qquad H_0=0.
\]
The rising factorial is $(a)_n=a(a+1)\cdots(a+n-1)$; the falling
factorial is $a^{\underline n}=a(a-1)\cdots(a-n+1)$.
The second-kind Stirling number is
$\left\{\begin{smallmatrix}n\\k\end{smallmatrix}\right\}$;
$s(n,k)$ denotes the signed first-kind Stirling number.
Coefficient extraction $[u^j]F$ always refers to an ordinary power
series, with any exponential-generating factorials displayed
separately.

For a function with a finite logarithmic singular expansion at zero,
\begin{equation}\label{scope:FP}
 \FP\int_0^1 f(x)\,dx
   =\CT_{\varepsilon\downarrow0}\int_\varepsilon^1 f(x)\,dx.
\end{equation}
The constant term means the coefficient of
$\varepsilon^0(\log\varepsilon)^0$ after discarding divergent powers
and logarithms. It is a coordinate-dependent prescription. A spectral
Laurent constant $\CT_{u=0}F(u)$ is a different operation; identifying
the two requires a proof of their local counterterms. Such a proof is
given in Section~\ref{cub:sec}.

For periodic distributions, $\delta$ is the unit Dirac mass at the
origin of $\mathbb R/\mathbb Z$, $D$ is the distribution derivative,
and $\check F(x)=F(-x)$. We normalize convolution so that
\[
 (\check F*G)(a)=\int_0^1 F(x)G(x+a)\,dx
\]
whenever ordinary integration is valid. A finite part of a pointwise
product is not defined by multiplying distributions. Section
\ref{cont:sec} explicitly compares two independently defined
extensions of the same off-collision function.

\paragraph{Why the directions matter.}
There is no single holomorphic Taylor germ of the uncolored Tornheim
function at $(0,0,0)$. For example,
\[
 T(t,t,t)\big|_{t=0}^{\mathrm{cont}}=\frac13,\qquad
 T(0,0,t)\big|_{t=0}^{\mathrm{cont}}=\frac5{12}.
\]
A notation such as $T''(0)$ without a specified one-variable
restriction is therefore insufficient. The same distinction occurs
when one exponent of a harmonic sum is fixed before differentiating
another. Every derivative formula below preserves its stated
restriction.

% END sections/01_scope.tex

% BEGIN sections/02_cubic.tex
\section{Coincident products: an explicit holomorphic completion}
\label{cub:sec}

The newest coincident-point report evaluates all quadratic Stieltjes
moments and explicitly asks for the cubic case, beginning with
$\FP\int_0^1\psi(x)^3\,dx$ \cite{IncomingCoincident}.
At three factors the local subtraction has a new feature: a pair of
singular factors and the linear term of the remaining factor together
produce a logarithmic divergence.  We give its exact spectral
counterterm.  The construction works at every derivative order and,
locally, at every number of factors.

\subsection{A finite local rule for any number of factors}

Use the lowering spectral convention
\begin{equation}\label{cub:generator}
 g(\lambda,x)=\zeta(1-\lambda,x)+\frac1\lambda
   =\sum_{m\ge0}\gamma_m(x)\frac{\lambda^m}{m!}.
\end{equation}
The pole at $\lambda=0$ is removable at each $x>0$.
For $r\ge0$ put
\[
 c_r(\lambda)=(-1)^r(1-\lambda)_r,
 \qquad h_r(\lambda,x)=\partial_x^r g(\lambda,1+x).
\]
The Hurwitz shift relation gives the exact splitting
\begin{equation}\label{cub:local-split}
 \partial_x^r g(\lambda,x)
 =c_r(\lambda)x^{\lambda-r-1}+h_r(\lambda,x).
\end{equation}
Every $h_r$ is holomorphic jointly in $\lambda$ near zero and $x$ near
zero.  Its required Taylor coefficients are explicit:
\begin{equation}\label{cub:h-coefficients}
 [x^k]h_r(\lambda,x)=
 \begin{cases}
 \zeta(1-\lambda)+1/\lambda,&r=k=0,\\[2pt]
 c_{r+k}(\lambda)\zeta(1+r+k-\lambda)/k!,&r+k>0.
 \end{cases}
\end{equation}

Let $d\ge2$ and $\boldsymbol r=(r_1,\ldots,r_d)$ be nonnegative
integers.  For $\varnothing\ne S\subseteq\{1,\ldots,d\}$ write
\[
 \lambda_S=\sum_{i\in S}\lambda_i,
 \qquad \nu_S=\sum_{i\in S}(r_i+1)-1,
\]
and define the finite numerator
\begin{equation}\label{cub:subset-numerator}
 A_S(\boldsymbol\lambda)=
 \left(\prod_{i\in S}c_{r_i}(\lambda_i)\right)
 [x^{\nu_S}]\prod_{j\notin S}h_{r_j}(\lambda_j,x).
\end{equation}
The empty smooth product is $1$.  In particular, $A_{\{1,\ldots,d\}}=0$
because $\nu_{\{1,\ldots,d\}}\ge1$.
Let $\mathcal L_{\boldsymbol r}$ denote the meromorphic continuation
of the ordinary integral of $\prod_i\partial_x^{r_i}g(\lambda_i,x)$.
It is initially defined, for example, with every
$\Re\lambda_i>r_i+1$, away from the removable parameter expressions.

\begin{theorem}[Subset completion in the unit coordinate]
\label{cub:general-completion}
The expression
\begin{equation}\label{cub:completion}
 \mathcal F_{\boldsymbol r}(\boldsymbol\lambda)
 =\mathcal L_{\boldsymbol r}(\boldsymbol\lambda)
   -\sum_{\varnothing\ne S\subsetneq\{1,\ldots,d\}}
            \frac{A_S(\boldsymbol\lambda)}{\lambda_S}
\end{equation}
is holomorphic in a neighborhood of the spectral origin.  Its Taylor
coefficients are exactly
\begin{equation}\label{cub:all-moments}
 \FP\int_0^1\prod_{i=1}^d\gamma_{m_i}^{(r_i)}(x)\,dx
 =\left(\prod_i m_i!\right)
 [\lambda_1^{m_1}\cdots\lambda_d^{m_d}]
          \mathcal F_{\boldsymbol r}(\boldsymbol\lambda).
\end{equation}
Here $\FP$ is the constant term of the cutoff integral in the fixed
coordinate $x$, with upper endpoint $1$ and no additional finite term.
The complete numerators in \eqref{cub:completion} are subtracted;
restricting a numerator to its polar hyperplane generally changes the
finite answer.
\end{theorem}

\begin{proof}
Fix $0<b<1$.  Expanding the exact product in
\eqref{cub:local-split} gives one term for each singular subset $S$,
namely
\[
 \left(\prod_{i\in S}c_{r_i}(\lambda_i)\right)
 x^{\lambda_S-\nu_S-1}
 \prod_{j\notin S}h_{r_j}(\lambda_j,x).
\]
For each nonempty $S$, subtract the Taylor polynomial of its smooth
product through degree $\nu_S$.  The remainder is normally integrable
on $(0,b)$ for all parameters in a sufficiently small common polydisc;
fixed parameter derivatives only add powers of $\log x$.
The Taylor monomial of degree $k$ has integral proportional to
\[
 \frac{b^{\lambda_S+k-\nu_S}}{\lambda_S+k-\nu_S}.
\]
Only $k=\nu_S$ has a denominator vanishing at the spectral origin.
Its meromorphic integral is $A_S b^{\lambda_S}/\lambda_S$.
The generating function of its coordinate finite parts is
$A_S(b^{\lambda_S}-1)/\lambda_S$.  Their difference is precisely
$A_S/\lambda_S$.  All other denominators are nonzero near the origin
and their Taylor coefficients equal the corresponding cutoff constant
terms.  The integral on $[b,1]$ is holomorphic and has the same
coefficient-extraction property.  Summing the finitely many subsets
proves both assertions.  Changing $b$ redistributes ordinary integrals
and Taylor integrals equally on both sides, so the result does not
depend on this auxiliary split.
\end{proof}

This theorem is a local completion rule, not a claim that an arbitrary
depth integral reduces to single zeta values.  For three factors its
global coordinates are the classical Tornheim function and ordinary
zeta functions.  That explicit global formula comes next.

\subsection{The cubic global kernel and all its argument derivatives}

Define, initially in its domain of absolute convergence,
\begin{equation}\label{cub:T-def}
 T(A,B,C)=\sum_{m,n\ge1}\frac{1}{m^A n^B(m+n)^C}.
\end{equation}
We always use its meromorphic continuation, with the order of any
restriction specified.  The classical Hurwitz--Tornheim connection is
due to earlier work, including Espinosa and Moll \cite{EM2006}; the
contribution here is its complete coordinate normalization at the
coincident Stieltjes point.

Put $\Lambda=\lambda_1+\lambda_2+\lambda_3$.  If $\{i,j,k\}$ is
$\{1,2,3\}$, define
\begin{align}
 K_2(a,b)&=\frac{2\Gamma(a)\Gamma(b)}{(2\pi)^{a+b}}
       \cos\frac{\pi(a-b)}2\,\zeta(a+b),\label{cub:K2}\\
 K_3(\boldsymbol\lambda)
 &=\frac{2\prod_i\Gamma(\lambda_i)}{(2\pi)^\Lambda}
   \sum_{k=1}^3
    \cos\frac{\pi(\lambda_i+\lambda_j-\lambda_k)}2
       T(\lambda_i,\lambda_j,\lambda_k).
       \label{cub:K3}
\end{align}
In \eqref{cub:K3}, the two indices other than $k$ may be taken in
increasing order, since $T(A,B,C)=T(B,A,C)$.

\begin{proposition}[Finite spectral expression]
\label{cub:global}
The kernels $K_2$ and $K_3$ are the meromorphic continuations of the
integrals of two and three factors $\zeta(1-\lambda_i,x)$,
respectively.  For $d=3$, the meromorphic integral in
Theorem~\ref{cub:general-completion} is explicitly
\begin{align}
 \mathcal L_{\boldsymbol r}(\boldsymbol\lambda)
 =\left(\prod_i c_{r_i}(\lambda_i)\right)
 \bigg\{&K_3(\lambda_1-r_1,\lambda_2-r_2,\lambda_3-r_3)
       \notag\\
 &+\sum_{k:r_k=0}\frac{
       K_2(\lambda_i-r_i,\lambda_j-r_j)}{\lambda_k}
       +\frac{\mathbf1_{\boldsymbol r=\boldsymbol0}}
                    {\lambda_1\lambda_2\lambda_3}\bigg\}.
 \label{cub:Lr}
\end{align}
Together with \eqref{cub:subset-numerator}, this is a finite
Tornheim--zeta coefficient formula for every cubic Stieltjes moment,
at arbitrary Stieltjes and argument-derivative orders.
\end{proposition}

\begin{proof}
For $\Re\lambda_i>1$, the nonzero Fourier coefficient of
$\zeta(1-\lambda_i,x)$ at $n$ is
\[
 \Gamma(\lambda_i)(2\pi|n|)^{-\lambda_i}
              e^{-\pi i\lambda_i\operatorname{sgn}(n)/2}.
\]
The Fourier series are absolutely convergent.  Integrating their
product keeps $n_1+n_2+n_3=0$, with all $n_i\ne0$.
There are three choices of the negative frequency and their three
sign reversals.  On each positive cone the negative magnitude is
$m+n$, giving exactly \eqref{cub:K3}.  The two-factor calculation
gives \eqref{cub:K2}; it is the classical Hurwitz product formula
\cite{EM2002}.  Finite endpoint Taylor subtraction supplies
meromorphic continuation of the original integrals, and uniqueness
extends these identities.  Finally,
$\partial_x^r g=c_r\zeta(1+r-\lambda,x)$ for $r>0$, whereas
$g=\zeta(1-\lambda,x)+1/\lambda$ for $r=0$.
Expand the three factors.  Integrals with just one Hurwitz factor
vanish by meromorphic continuation of
$\int_0^1\zeta(s,x)\,dx=0$.  The other terms give \eqref{cub:Lr}.
\end{proof}

For undifferentiated factors set
$h(\lambda)=\zeta(1-\lambda)+1/\lambda$.  Formula
\eqref{cub:completion} becomes particularly short:
\begin{align}
 \mathcal F_{\boldsymbol0}={}&K_3+
   \sum_{k=1}^3\frac{K_2(\lambda_i,\lambda_j)}{\lambda_k}
   +\frac{1}{\lambda_1\lambda_2\lambda_3}
   -\sum_{i=1}^3\frac{h(\lambda_j)h(\lambda_k)}{\lambda_i}
       \notag\\
 &+\sum_{k=1}^3\frac{(1-\lambda_k)\zeta(2-\lambda_k)}
                              {\lambda_i+\lambda_j}.
 \label{cub:Fexplicit}
\end{align}
The last line is necessary.  It is the pair-subset correction absent
from the separated-singularity triple formula.  Neither it nor the
singleton correction can be suppressed before restricting spectral
directions.

\subsection{The stated cubic digamma target}

Let
\[
 \omega(t)=T(t,t,t),\quad L=\log(2\pi),\quad
 E(t)=\Gamma(1+t)e^{-Lt},\quad
 \gamma=\gamma_0(1).
\]
The symbol $\omega^{(j)}(0)$ means a derivative of this specified
one-variable continuation.  It does not mean a derivative of a
jointly holomorphic $T$ at $(0,0,0)$.

\begin{theorem}[Diagonal second derivative and cubic digamma moment]
\label{cub:diagonal}
The diagonal continuation is holomorphic at zero and satisfies
\begin{equation}\label{cub:omega-jets}
 \omega(0)=\frac13,\qquad \omega'(0)=L,
 \qquad \boxed{\omega''(0)=L^2-4\zeta''(0)}.
\end{equation}
The first two values are classical \cite{BorweinDilcher}.
The coordinate-normalized cubic moment is
\begin{align}
 \boxed{\displaystyle\FP\int_0^1\psi(x)^3\,dx}
 ={}&-\omega'''(0)-8\zeta'''(0)
       -12(\gamma+L)\zeta''(0)-L^3-6\gamma L^2
       \notag\\
 &+5\gamma^3+6\gamma\gamma_1
       -\frac32\gamma\zeta(2)
       +\frac32\bigl(\zeta(2)+\zeta'(2)\bigr).
 \label{cub:psi-cube}
\end{align}
Thus the requested cubic finite part is exactly reduced to one third
derivative of the diagonal Tornheim function and ordinary zeta and
Stieltjes constants.  No reduction of $\omega'''(0)$ to ordinary zeta
constants is asserted.
\end{theorem}

\begin{proof}
On the diagonal, \eqref{cub:Fexplicit} reads
\begin{align}
 F(t):=\mathcal F_{\boldsymbol0}(t,t,t)
 ={}&\frac{6E(t)^3\cos(\pi t/2)\omega(t)
              +6E(t)^2\zeta(2t)+1}{t^3}
       \notag\\
 &-\frac{3h(t)^2}{t}
       +\frac{3(1-t)\zeta(2-t)}{2t}.
 \label{cub:Fdiagonal}
\end{align}
The left side is holomorphic by
Theorem~\ref{cub:general-completion}.  Solving this identity for
$\omega(t)$ proves its holomorphy at zero, since $E(0)=1$ and
$\cos0=1$.  Write
\[
 E(t)=\exp\left[-(\gamma+L)t+
                 \frac{\zeta(2)}2t^2-\frac{\zeta(3)}3t^3+O(t^4)\right].
\]
Cancellation of the coefficients of $t^{-3}$ and $t^{-2}$ gives the
first two values of \eqref{cub:omega-jets}, using
$\zeta(0)=-1/2$ and $\zeta'(0)=-L/2$.
The numerator coefficient at $t^2$ is
\[
 3\omega''(0)+3\gamma^2-3L^2-\tfrac32\zeta(2)+12\zeta''(0).
\]
The other two terms of \eqref{cub:Fdiagonal} contribute
$-3\gamma^2+3\zeta(2)/2$ to the coefficient of $t^{-1}$.
This proves the second-derivative formula.

The numerator coefficient at $t^3$, after those substitutions, is
\[
 \omega'''(0)+8\zeta'''(0)+12(\gamma+L)\zeta''(0)
 +L^3+6\gamma L^2-5\gamma^3+\tfrac32\gamma\zeta(2).
\]
The remaining constant terms are
$-6\gamma\gamma_1-3(\zeta(2)+\zeta'(2))/2$.
Finally $F(0)=\FP\int\gamma_0(x)^3dx=-\FP\int\psi(x)^3dx$,
which proves \eqref{cub:psi-cube}.
\end{proof}

A direct, convergent integral gives an independent way to evaluate
the same left side.  The expansion
$\psi(x)=-x^{-1}-\gamma+\zeta(2)x-\zeta(3)x^2+\cdots$ gives
\begin{align}
 \FP\int_0^1\psi(x)^3\,dx
 ={}&\int_0^1\left[\psi(x)^3+\frac1{x^3}
       +\frac{3\gamma}{x^2}
       -\frac{3\zeta(2)-3\gamma^2}{x}\right]dx
       +\frac12+3\gamma.
 \label{cub:direct-cube}
\end{align}
The integrand has a finite limit at zero.  Independent truncated
Mellin-series evaluation of $\omega'''(0)$ and local-series quadrature
of \eqref{cub:direct-cube} agree to more than thirty decimal places:
\begin{align}
 \omega''(0)&=11.403217934867761113621595961662\ldots,\notag\\
 \omega'''(0)&=86.657866601100073652903276858762\ldots,\notag\\
 \FP\int_0^1\psi(x)^3dx
   &=1.966262734391534340664375338411\ldots.
 \label{cub:numerics}
\end{align}
These are numerical diagnostics accompanying the analytic proof;
the implementation does not claim interval-certified final digits.

\begin{corollary}[A cubic derivative specialization]
\label{cub:cubic-derivative}
In the same coordinate convention,
\begin{equation}\label{cub:psi-square-prime}
 \FP\int_0^1\psi(x)^2\psi'(x)\,dx
             =\zeta(3)-2\gamma\zeta(2).
\end{equation}
\end{corollary}
\begin{proof}
Take the constant term of
$\int_\epsilon^1\psi^2\psi'=(\psi(1)^3-\psi(\epsilon)^3)/3$.
The constant coefficient of $\psi(\epsilon)^3$ is
$-\gamma^3+6\gamma\zeta(2)-3\zeta(3)$.
\end{proof}

More generally, a finite-part integral of a derivative of a product
equals its upper-endpoint value minus the constant term of its local
expansion.  This supplies a terminating family of relations between
the cubic moments in \eqref{cub:all-moments}; it does not impose a
spurious independent value on the remaining primitive cubic moment.

\subsection{Why a spectral direction must be recorded}

The reduction at fixed $A=B=0$ gives
\begin{equation}\label{cub:C-axis}
 T(0,0,C)=\zeta(C-1)-\zeta(C).
\end{equation}
Indeed, $m+n=N$ has $N-1$ positive solutions.  Continuing this
one-variable identity gives $5/12$ at $C=0$, whereas
\eqref{cub:omega-jets} gives $1/3$ on the diagonal.
This direction dependence is already explicit in
Borwein and Dilcher \cite[Theorem~2.7]{BorweinDilcher}.
It is not an inconsistency.  The multivariable function is meromorphic
at intersecting polar hyperplanes, and its restrictions carry different
finite information.  The completed function
$\mathcal F_{\boldsymbol r}$ is genuinely holomorphic; the uncompleted
$T$ generally is not.  Section~\ref{sp:sec:rays} computes the first two
derivatives along an arbitrary transverse ray.

% END sections/02_cubic.tex

% BEGIN sections/03_contacts.tex
\section{The full distributional collision correction}
\label{cont:sec}

The incoming coincident-point report isolates the off-point collision
singularity and its finite constant, but leaves the distribution
supported at the collision undetermined \cite{IncomingCoincident}.
We now specify the extension convention and calculate that distribution
for every pair of Stieltjes indices and every pair of argument
derivatives.  This is a comparison of two extensions of the same
off-point function.  It does not identify a pointwise product of two
singular distributions at the same argument.

\subsection{The two extensions being compared}

Work on $\mathbb T=\mathbb R/\mathbb Z$.  Let $G_m$ be the unit-coordinate
Hadamard extension of $\gamma_m(x)$ from $0<x<1$, and let
$F_{m,p}$ be the corresponding extension of
$\gamma_m^{(p)}(x)$.  Thus, for a smooth periodic test function,
\begin{equation}\label{cont:finite-part}
 \langle F_{m,p},\phi\rangle
 =\CT_{\epsilon\downarrow0}
      \int_\epsilon^1\gamma_m^{(p)}(x)\phi(x)\,dx,
 \qquad F_{m,0}=G_m.
\end{equation}
Write $D=d/dx$, $\Delta=\delta_0-1$, and
$\check U(x)=U(-x)$.  The exact correlation is the well-defined
periodic convolution
\begin{equation}\label{cont:correlation}
 \mathcal R_{mn}^{p,q}=\check F_{m,p}*F_{n,q}.
\end{equation}
For ordinary smooth functions this is
$\int_0^1 f(x)g(x+a)\,dx$ as a function of $a$.

For $0<a<1$, let $J_{mn}^{p,q}(a)$ be the unit-coordinate finite part
of the separated product
$\int_0^1\gamma_m^{(p)}(x)\gamma_n^{(q)}(\{x+a\})\,dx$.
It is the off-point restriction of \eqref{cont:correlation}, as
established in the manuscript's translated contact law
\cite{Repo}.  Define its canonical extension in the shift coordinate by
\begin{equation}\label{cont:shift-extension}
 \left\langle\mathfrak F J_{mn}^{p,q},\phi\right\rangle
 =\CT_{\epsilon,\eta\downarrow0}
      \int_\epsilon^{1-\eta}J_{mn}^{p,q}(a)\phi(a)\,da.
\end{equation}
The two endpoint constants are taken independently in $a$ and $1-a$;
there is no added delta term.  The local power--log expansions proved
in the incoming collision report guarantee existence.

\subsection{Finite generating polynomials}

This section uses a raising parameter to simplify the Gamma quotient:
\[
 \mathcal G(u)=\sum_{m\ge0}(-1)^mG_m\frac{u^m}{m!},
 \qquad \mathcal F_p(u)=\sum_{m\ge0}(-1)^mF_{m,p}\frac{u^m}{m!}.
\]
For $r\ge0$ set
\begin{equation}\label{cont:P-H}
 P_r(z)=\prod_{j=1}^r(1+z/j),\qquad
 H_r(z)=\frac{P_r(z)-1}{z},\qquad P_0=1,\ H_0=0.
\end{equation}
$H_r$ is a polynomial, not the scalar harmonic number; its constant
coefficient is $\sum_{j=1}^r1/j$.  All its coefficients are elementary
symmetric polynomials in $1,1/2,\ldots,1/r$.

Put $w=u+v$ and define
\begin{equation}\label{cont:A-B}
 A(u,v)=\frac{\Gamma(1-u)\Gamma(1+w)}{\Gamma(1+v)},
 \qquad T_0(u,v)=\frac{A(u,v)-1}{uw},
 \qquad B(u,v)=T_0(u,v)+T_0(v,u).
\end{equation}
The quotients are removable because $A(0,v)=A(u,-u)=1$.
The convergent local series
\begin{equation}\label{cont:log-A}
 \log A(u,v)=\sum_{k\ge2}\frac{\zeta(k)}{k}
       \left[u^k+(-1)^k\big((u+v)^k-v^k\big)\right]
\end{equation}
makes every coefficient an exact finite polynomial in positive
integer zeta values.  The undifferentiated off-point use of this
Gamma quotient is inherited from the manuscript \cite{Repo}.

There is also a compact symmetric form:
\begin{equation}\label{cont:B-symmetric}
 B(u,v)=\frac1{uv}\left\{
  \frac{\Gamma(1-u)\Gamma(1-v)}{\Gamma(1-u-v)}
  \frac{\cos(\pi(u-v)/2)}{\cos(\pi(u+v)/2)}-1\right\}.
\end{equation}
Indeed,
$uvB=(vA(u,v)+uA(v,u))/(u+v)-1$.
The Gamma reflection formula turns the fraction before the subtraction
into the Gamma quotient in \eqref{cont:B-symmetric} times
$\{\sin(\pi u)+\sin(\pi v)\}/\sin(\pi(u+v))$.
The sine addition formula gives the displayed cosine ratio. All
apparent quotients are interpreted by their removable holomorphic
extensions near $(0,0)$. This second form provides an independent
coefficient check on the generator.

\begin{theorem}[All-index collision contact distribution]
\label{cont:main}
Let $m,n,p,q\ge0$ and $r=p+q$.  Then
\begin{equation}\label{cont:main-identity}
 \boxed{\mathcal R_{mn}^{p,q}
      =\mathfrak F J_{mn}^{p,q}+e_{mn}^{p,q}\delta_0^{(r)}}.
\end{equation}
There are no additional lower derivatives of $\delta_0$.
The coefficient is explicitly
\begin{equation}\label{cont:coefficient}
 e_{mn}^{p,q}=(-1)^{m+n}m!n![u^mv^n]E_{p,q}(u,v),
\end{equation}
where the following holomorphic germ requires only polynomial
arithmetic and \eqref{cont:log-A}:
\begin{align}
 E_{p,q}(u,v)=(-1)^p\bigg\{
 &P_r(w)B(u,v)
   +\frac{H_r(w)-H_r(v)}{u}
   +\frac{H_r(w)-H_r(u)}{v}\notag\\
 &-H_p(u)H_r(v)-H_q(v)H_r(u)+H_p(u)H_q(v)\bigg\}.
 \label{cont:E}
\end{align}
Both divided differences in this formula are polynomials.
\end{theorem}

\begin{proof}
Let $Z_s$ be the meromorphic periodic Hurwitz distribution.  Its
Fourier coefficients are the Hurwitz coefficients used in
Section~\ref{cub:sec}, with zero mean.  Local Taylor subtraction at
the origin gives
\begin{equation}\label{cont:Z-expansion}
 Z_{1+u}=-\frac\Delta u+\mathcal G(u).
\end{equation}
For $p>0$, the same subtraction at the higher Hurwitz pole, together
with $D^pZ_s=(-1)^p(s)_pZ_{s+p}$, gives
\begin{equation}\label{cont:derivative-anomaly}
 \mathcal F_p(u)=D^p\mathcal G(u)+H_p(u)\delta_0^{(p)}.
\end{equation}
This generating version of the manuscript's differentiation anomaly
also holds for $p=0$.  For clarity, the polar term of
$Z_{1+p+u}$ is $(-1)^{p+1}\delta_0^{(p)}/(p!u)$.
Multiplication by $(-1)^p(1+u)_p$ produces
$-P_p(u)\delta_0^{(p)}/u$; subtracting the pole of the left side
leaves exactly \eqref{cont:derivative-anomaly}.

The beta form of the classical correlation, continued as a periodic
distribution, is
\[
 \check Z_{1+u}*Z_{1+v}
 =-\frac{A(u,v)}u Z_{1+w}
   -\frac{A(v,u)}v\check Z_{1+w}.
\]
It follows either from the Fourier coefficients or first in a
convergent half-plane and then by meromorphic continuation.
Substitute \eqref{cont:Z-expansion} and use
$\Delta*\Delta=\Delta$ and $\Delta*G_m=G_m$.
The result is the holomorphic distribution identity
\begin{align}
 \check{\mathcal G}(u)*\mathcal G(v)
 ={}&B(u,v)\Delta
 -\frac{\mathcal G(w)-\mathcal G(v)}u
 -\frac{\check{\mathcal G}(w)-\check{\mathcal G}(u)}v
       \notag\\
 &-w\bigl[T_0(u,v)\mathcal G(w)
              +T_0(v,u)\check{\mathcal G}(w)\bigr].
 \label{cont:base-generator}
\end{align}
Its off-point constant is $-B$.  The canonical shift extension
replaces each ordinary Stieltjes function by its corresponding $G_m$
and retains that constant.  Consequently the difference for $r=0$
is $B\delta_0$, as claimed; it is not $B\Delta$.

For $r>0$, converting $D^r\mathcal G(z)$ to the canonical extension
of its ordinary derivative adds the difference
$-H_r(z)\delta_0^{(r)}$.  A reflected term has the same difference:
reflection and the chain rule each contribute $(-1)^r$.
Applying this to \eqref{cont:base-generator}, the contact coefficient
of its $r$th derivative is
\[
 P_r(w)B+
 \frac{H_r(w)-H_r(v)}u+
 \frac{H_r(w)-H_r(u)}v.
\]
Finally expand the correlation of
$D^p\mathcal G(u)+H_p(u)\delta_0^{(p)}$ and
$D^q\mathcal G(v)+H_q(v)\delta_0^{(q)}$.
The two cross terms and the delta--delta term supply, respectively,
$-H_p(u)H_r(v)$, $-H_q(v)H_r(u)$, and $H_p(u)H_q(v)$.
The orientation of the first factor contributes $(-1)^p$.
This proves \eqref{cont:E}.  Every contact term in this calculation
has derivative order $r$, proving the absence of lower orders.
Coefficient extraction gives \eqref{cont:main-identity}.
\end{proof}

\subsection{Explicit contact coefficients}

For undifferentiated factors, $E_{0,0}=B$ and the first values are
\begin{center}
\begin{tabular}{c c c}
\toprule
$(m,n)$ & $e_{mn}^{0,0}$ & contact term\\
\midrule
$(0,0)$ & $2\zeta(2)$ & $2\zeta(2)\delta_0$\\
$(1,0)$ or $(0,1)$ & $\zeta(3)$ & $\zeta(3)\delta_0$\\
$(2,0)$ or $(0,2)$ & $11\zeta(4)/2$ & $11\zeta(4)\delta_0/2$\\
$(1,1)$ & $7\zeta(4)/2$ & $7\zeta(4)\delta_0/2$\\
\bottomrule
\end{tabular}
\end{center}
The last two simplifications use $\zeta(2)^2=5\zeta(4)/2$.
For Stieltjes indices $m=n=0$ and positive argument derivatives,
\begin{align}
 e_{00}^{1,0}&=1-2\zeta(2),\notag\\
 e_{00}^{1,1}&=1-2\zeta(2),\notag\\
 e_{00}^{2,0}&=2\zeta(2)-\frac54.
 \label{cont:low-examples}
\end{align}
The associated distributions are $\delta_0'$, $\delta_0''$, and
$\delta_0''$, respectively.  Swapping the two factors reflects the
correlation; accordingly
\begin{equation}\label{cont:reflection}
 e_{nm}^{q,p}=(-1)^r e_{mn}^{p,q}.
\end{equation}
This identity is also immediate from \eqref{cont:E}.

For instance, the known off-point identity is
$J_{00}^{0,0}(a)=\gamma_1(a)+\gamma_1(1-a)-2\zeta(2)$.
Its exact periodic distributional completion is
\begin{equation}\label{cont:example00}
 \check G_0*G_0=G_1+\check G_1+2\zeta(2)(\delta_0-1).
\end{equation}
An off-point numerical computation cannot see the delta term.
In contrast, a nonzero Fourier coefficient of \eqref{cont:example00}
contains the necessary constant $2\zeta(2)$; omitting it gives the
wrong correlation even though the two functions agree for $0<a<1$.

\subsection{Normalized primitives of the contact discrepancy}

The contact formula also fixes constants in antiderivatives.
Let $\mathcal I$ be the zero-mean inverse of $D$ on periodic
distributions of mean zero.  For the periodic Bernoulli polynomial
$B_j(\{a\})$,
\begin{equation}\label{cont:Bernoulli-inverse}
 \mathcal I^j\Delta=-\frac{B_j(\{a\})}{j!},\qquad j\ge1.
\end{equation}
This follows from $DB_1(\{a\})=1-\delta_0$ and
$B_j'=jB_{j-1}$ for $j>1$, with zero means.
For $r>0$, applying $\mathcal I^{r+j}$ to the discrepancy in
\eqref{cont:main-identity} yields
\begin{equation}\label{cont:contact-primitive}
 \mathcal I^{r+j}\bigl(\mathcal R_{mn}^{p,q}
             -\mathfrak FJ_{mn}^{p,q}\bigr)
       =-e_{mn}^{p,q}\frac{B_j(\{a\})}{j!},\qquad j\ge1.
\end{equation}
For $r=0$, first subtract the mean of the discrepancy, replacing
$e\delta_0$ by $e\Delta$, and the same formula applies.
These are fixed normalized primitives, with no unspecified additive
polynomial.  They explain exactly how a point-supported correction
can become an ordinary Bernoulli term after repeated integration.

% END sections/03_contacts.tex

% BEGIN sections/04_coordinates.tex
% Integrable regular parts may be dropped in all local coordinate comparisons.
% Requires amsmath, amssymb, amsthm.  Labels are prefixed invg:.
% Suggested bibliography keys: IncomingExact, IncomingTransport, Latt2011.
\section{Which logarithmic singularities have no coordinate defect?}
\label{invg:sec}

The nonlinear coordinate formulas in the two incoming reports give a
finite algorithm for each specified singular germ.  The fifth research
question in the Exact Identities report asks for an intrinsic
classification of the germs whose defect vanishes under every coordinate
tangent to the identity.  We answer that question, and more generally
allow any prescribed order of tangency.  The necessary condition is as
important as the sufficient one: logarithmic degrees from different pole
orders cannot conceal a defect for every coordinate.

The change-of-variables mechanism for logarithmic finite parts is
classical; see, for example, L\"att's Theorem~2 in
\cite{Latt2011}.  The theorem below is an exact invariant-subspace
classification and a specialization to the Stieltjes hierarchy relative
to the audited incoming material.  No general historical-priority claim
is made.

\subsection{Conventions and the inherited residue formula}

Work locally on the positive side of the origin.  A logarithmic
meromorphic germ has the unique singular expansion
\begin{equation}\label{invg:germ}
 h(x)=\sum_{q=1}^{R}\sum_{n=0}^{L}c_{q,n}x^{-q}\log^n x+h_{\rm reg}(x),
\end{equation}
where the omitted part is locally integrable and is a finite sum of
analytic germs times nonnegative powers of $\log x$.  Coefficients may be
real or complex.  Define $\operatorname{FP}_x h$ by taking the constant
term of its cutoff integral against a smooth test function.  Positive
powers of the cutoff logarithm and negative powers of the cutoff are
removed.  Use scalar distribution pullback, so that ordinary functions
pull back by composition.

For an orientation-preserving analytic coordinate $f$, put
\[
 H_f(t)=f(t)/t,\qquad \lambda_f(t)=\log H_f(t),\qquad
 C_f[h]=f^*(\operatorname{FP}_x h)-\operatorname{FP}_t(h\circ f).
\]
The residue law proved in the incoming reports is
\begin{equation}\label{invg:residue}
 \langle C_f[h],\phi\rangle
 =\sum_{q,n}\frac{c_{q,n}}{n+1}
    \operatorname*{Res}_{t=0}
     f(t)^{-q}\lambda_f(t)^{n+1}\phi(t)\,dt.
\end{equation}
It is enough to keep the finite Taylor jet of the smooth test function
inside the residue.  Integrable terms create no discrepancy.  Equivalently,
for a single monomial,
\begin{equation}\label{invg:monomial}
 C_f[x^{-q}\log^n x]
 =n!\sum_{j=0}^{q-1}\frac{(-1)^j}{j!}
 [z^{n+1}t^{q-1-j}]H_f(t)^{z-q}\,\delta_0^{(j)}.
\end{equation}
This is the same coefficient law, because
$[z^{n+1}]H_f^{z-q}=H_f^{-q}\lambda_f^{n+1}/(n+1)!$.

For completeness, the analytic justification is short.  Continue
$x_+^{z-q}$ after subtracting the test-function Taylor polynomial through
degree $q-1$.  Only the term of degree $q-1$ has denominator $z$ near
zero.  After pullback, its numerator depends on $z$ through $H_f^{z-q}$.
The constant term of the cutoff coefficient omits the quantity
$n![z^{n+1}]$ of that numerator, whereas the meromorphic coefficient
retains it.  Every other Taylor denominator is nonzero at $z=0$; the
remainder is normally integrable there.  This proves
\eqref{invg:monomial}, including its sign and factorials.

\subsection{The complete classification}

For an integer $d\ge1$, write
\[
 \mathscr G_d=\{f:f(0)=0,\quad f(t)=t+O(t^{d+1})\}.
\]
These are germs with identity Taylor jet through degree $d$.

\begin{theorem}[Classification by pole order and logarithmic degree]
\label{invg:classification}
For the germ \eqref{invg:germ} and an integer $d\ge1$, the following
conditions are equivalent.
\begin{enumerate}
\item $C_f[h]=0$ for every $f\in\mathscr G_d$.
\item $C_{f_a}[h]=0$ for every sufficiently small real $a$, where
      $f_a(t)=t(1+a t^d)$.
\item $c_{q,n}=0$ whenever $q>d(n+1)$.
\end{enumerate}
Equivalently, if $h=\sum_{n=0}^{L}a_n(x)\log^n x$, the pole order of
$a_n$ is at most $d(n+1)$ for every $n$.  For $d=1$, the exact condition
is $\operatorname{poleord}(a_n)\le n+1$.
\end{theorem}

\begin{proof}
The first condition implies the second.  To prove that the third implies
the first, note that $\lambda_f(t)=O(t^d)$ for $f\in\mathscr G_d$.
The $q,n$ summand of the residue integrand in \eqref{invg:residue} has
order at least $d(n+1)-q\ge0$.  It is analytic at zero, so its residue
against any smooth test function is zero.

It remains to prove necessity; this is not a consequence of estimating
the order of one monomial alone.  For $f_a=t(1+a t^d)$,
\eqref{invg:monomial} becomes
\begin{equation}\label{invg:quadratic-family}
 C_{f_a}[x^{-q}\log^n x]
 =n!\sum_{k=1}^{\lfloor(q-1)/d\rfloor}
 \frac{(-1)^{q-1-dk}a^k}{(q-1-dk)!}
 [z^{n+1}]\binom{z-q}{k}\,\delta_0^{(q-1-dk)}.
\end{equation}
The term at $k=0$ vanishes because $n+1\ge1$.
For a fixed power $a^k$ and derivative $\delta_0^{(j)}$, the pole order
is uniquely determined by $q=dk+j+1$.  Therefore different pole orders
cannot cancel this coefficient.

Fix $q$ and let $K=\lfloor(q-1)/d\rfloor$.  The coefficient of
$a^k\delta_0^{(q-1-dk)}$ vanishes for $k=1,\ldots,K$ only if
\begin{equation}\label{invg:triangular}
 \sum_{n=0}^{k-1}c_{q,n}n![z^{n+1}]\binom{z-q}{k}=0,
 \qquad k=1,\ldots,K.
\end{equation}
This is triangular in $c_{q,0},\ldots,c_{q,K-1}$, and its diagonal
coefficient at $c_{q,k-1}$ is
\[
 (k-1)![z^k]\binom{z-q}{k}=\frac1k\ne0.
\]
Induction on $k$ gives $c_{q,0}=\cdots=c_{q,K-1}=0$.
Those indices are exactly the pairs satisfying $q>d(n+1)$.
This proves the third condition and completes the equivalence.
\end{proof}

A useful consequence is that the least required tangency order can be
read directly from the singular expansion:
\begin{equation}\label{invg:least-tangency}
 d_{\min}(h)=\max_{c_{q,n}\ne0}
              \left\lceil\frac q{n+1}\right\rceil.
\end{equation}
For a germ with no singular term, every coordinate is harmless.  For a
nonzero singular germ, \eqref{invg:least-tangency} is the least positive
integer $d$ such that matching the identity jet through degree $d$
removes every coordinate defect.

\begin{corollary}[Exact dimension of the obstruction space]
\label{invg:rank}
Within pole order at most $R$ and logarithmic degree at most $L$, the
codimension of the invariant subspace is
\begin{equation}\label{invg:rank-formula}
 \sum_{q=1}^{R}\min\!\left(L+1,\left\lfloor\frac{q-1}{d}\right\rfloor\right).
\end{equation}
The one-parameter family in \eqref{invg:quadratic-family}, with its
polynomial dependence on $a$ retained, detects all these obstructions.
\end{corollary}
\begin{proof}
Theorem~\ref{invg:classification} sets precisely the listed independent
coordinates $c_{q,n}$ to zero.  Equation~\eqref{invg:triangular} proves
that their obstruction functionals are independent.
\end{proof}

The reference to the whole parameter family matters: a single numerical
choice of $a$ can combine coefficients belonging to different pole orders.
The theorem uses an exact polynomial identity in $a$.

\subsection{Sharp Stieltjes thresholds for any order of tangency}

Use the usual convention
$\gamma_n(x)=x^{-1}\log^n x+\gamma_n(1+x)$.  Let
$T_{n,r}=\operatorname{FP}_x\gamma_n^{(r)}(x)$ be the finite part of
the pointwise argument derivative.  This object is distinguished from
$D^r(\operatorname{FP}_x\gamma_n)$, whose harmonic contact term was
computed in the incoming reports.

\begin{theorem}[Exact threshold for Stieltjes argument derivatives]
\label{invg:stieltjes-threshold}
For $r,n\ge0$ and $d\ge1$,
\[
 f^*T_{n,r}=\operatorname{FP}_t\bigl(\gamma_n^{(r)}\circ f\bigr)
 \quad\text{for every }f\in\mathscr G_d
\]
holds if and only if
\begin{equation}\label{invg:stieltjes-criterion}
 d\ge r+1\qquad\text{or}\qquad
 n\ge r+\left\lfloor\frac rd\right\rfloor.
\end{equation}
When $1\le d\le r$, set $K=\lfloor r/d\rfloor$, $e=r-dK$, and
$n_*=r+K-1$.  The first excluded index has the exact defect
\begin{equation}\label{invg:stieltjes-top}
 C_{f_a}[\gamma_{n_*}^{(r)}]
 =\frac{(-1)^e n_*!}{e!K!}a^K\delta_0^{(e)},
 \qquad f_a(t)=t(1+a t^d).
\end{equation}
Thus every threshold in \eqref{invg:stieltjes-criterion} is attained.
\end{theorem}

\begin{proof}
The singular part is
\begin{equation}\label{invg:stieltjes-singular}
 \gamma_n^{(r)}(x)
 =(-1)^rr!x^{-r-1}
 n![z^n]\left\{\prod_{j=1}^{r}\left(1-\frac zj\right)e^{z\log x}\right\}
 +\gamma_n^{(r)}(1+x).
\end{equation}
Every elementary symmetric coefficient of
$1,1/2,\ldots,1/r$ is nonzero.  Hence its smallest logarithmic degree is
$\max(n-r,0)$, and the coefficient at that degree is nonzero.
The classification theorem gives the exact condition
\[
 r+1\le d\bigl(\max(n-r,0)+1\bigr).
\]
If $d\ge r+1$ this holds for all $n$.  Otherwise it requires $n\ge r$
and is equivalent to the second inequality in
\eqref{invg:stieltjes-criterion}.

At $n=n_*$, the smallest logarithmic degree is $K-1$, with coefficient
$n_*!/(K-1)!$ in \eqref{invg:stieltjes-singular}.  Every higher
logarithmic degree is invariant.  In
\eqref{invg:quadratic-family}, only $k=K$ survives for this smallest
degree, and its derivative index is $e$.  Its coefficient is
$(-1)^ea^K/(e!K)$.  Multiplication by $n_*!/(K-1)!$ proves
\eqref{invg:stieltjes-top}.
\end{proof}

For $d=1$ this recovers the incoming sharp threshold $n\ge2r$ and its
coefficient $(2r-1)!a^r\delta_0/r!$.  It additionally gives, for example,
\[
 C_{t+a t^3}[\gamma_3^{(3)}]=-6a\delta_0',
 \qquad
 C_{t+a t^3}[\gamma_5^{(4)}]=60a^2\delta_0,
\]
while $\gamma_4^{(3)}$ and $\gamma_6^{(4)}$ have no defect under any
coordinate $f(t)=t+O(t^3)$.

\subsection{Consequences for coincident Stieltjes products}

The invariant space is linear but is generally not closed under
pointwise multiplication: $x^{-1}$ is invariant under every unit-tangent
coordinate, whereas $x^{-2}$ is not.  The classification nevertheless
provides a useful exact product criterion.

\begin{corollary}[Products with at most one digamma factor]
\label{invg:stieltjes-products}
Let $m_1,\ldots,m_s$ be nonnegative integers, at most one of which is
zero.  Then the specified pointwise-product finite part satisfies
\begin{equation}\label{invg:stieltjes-products-formula}
 f^*\operatorname{FP}_x\!\left(\prod_{j=1}^{s}\gamma_{m_j}(x)\right)
 =\operatorname{FP}_t\!\left(\prod_{j=1}^{s}\gamma_{m_j}(f(t))\right)
 \qquad (f\in\mathscr G_1).
\end{equation}
In particular, all powers $\operatorname{FP}_x\gamma_m(x)^s$ with
$m\ge1$ have this invariance.
\end{corollary}
\begin{proof}
Choose the singular term from a nonempty subset $S$ of the factors.
Before expanding the remaining analytic factors, the resulting term is
\[
 x^{-|S|}\log^{\sum_{j\in S}m_j}x
             \prod_{j\notin S}\gamma_{m_j}(1+x).
\]
Its logarithmic degree is at least $|S|-1$.  Taylor expansion of the
analytic factor can only decrease the pole order.  Thus every resulting
singular monomial satisfies $q\le n+1$, which is exactly the criterion
of Theorem~\ref{invg:classification} for $d=1$.
\end{proof}

The restriction on the number of zero indices cannot be replaced by a
large total Stieltjes index.  If $f(t)=t+a t^2+O(t^3)$ and
$\gamma_m=\gamma_m(1)$, direct use of the same classification gives
\begin{align}
 C_f[\gamma_0^2]&=a\delta_0,\label{invg:product-counter1}\\
 C_f[\gamma_0^2\gamma_1]&=
             \left(a\gamma_1+\frac{a^2}{2}\right)\delta_0,
             \label{invg:product-counter2}\\
 C_f[\gamma_0^2\gamma_m]&=a\gamma_m\delta_0\qquad(m\ge2).
             \label{invg:product-counter3}
\end{align}
Indeed, the mixed singular--singular--regular term in the last two
products is $\gamma_m x^{-2}$.  The fully singular term is
$x^{-3}\log^m x$: it contributes $a^2\delta_0/2$ when $m=1$ and
nothing when $m\ge2$.  All other terms meet the invariant criterion.
These are identities for a specified finite part of the ordinary
product; they do not postulate a product of distributions.

\paragraph{Verification.}
The accompanying exact script compares the residue construction with
\eqref{invg:quadratic-family}, computes the rank of the finite
obstruction map, and checks the Stieltjes criterion and its sharp
coefficient through several pole and tangency orders.  These are
independent algebraic audits of signs and finite formulas.  The
all-order conclusions follow from the proofs above.

% END sections/04_coordinates.tex

% BEGIN sections/05_spectral.tex
\section{Exact normal spectral jets at integer inner orders}
\label{sp:sec:normal}

Write
\[
 \mathcal T(A,B,C;X,Y)
 =\sum_{m,n\geq1}\frac{X^mY^n}{m^A n^B(m+n)^C}.
\]
The orders are independent complex variables, and powers of positive
integers use real logarithms.  The question in the incoming
\emph{Resonance and Coordinate Transport} report asks for finite
reductions of the normal $C$-derivatives at $C=0$, beginning with
$A,B\in\{0,1\}$.  The negative-$C$ value identity by itself does not
answer that question.  The reduction below keeps $C$ variable from
the outset.

The underlying positive-integer partial-fraction identity is classical:
Zhao explicitly uses it to reduce colored Tornheim values to double
polylogarithms \cite{sp:Zhao}.  Our use of that identity retains the
complete complex $C$-dependence, specifies the continuation required
by the normal jets, and combines it with a rational-function theorem
covering arbitrary depth, integer weights, and coincident colors.
These are exact reductions to stated function classes.  No arithmetic
independence or global priority claim is inferred from them.

For the double polylogarithm use the convention
\begin{equation}\label{sp:doubleLi}
 \operatorname{Li}_{u,v}(z,w)
   =\sum_{k>n\geq1}\frac{z^k w^n}{k^u n^v}.
\end{equation}
Initially the colors can be taken with $|X|,|Y|<1$ and $X,Y\ne0$,
so that the reindexed series below converge locally normally for all
orders.  On $|X|=|Y|=1$ we assume $X,Y\ne1$ and use radial analytic
continuation: replace $(X,Y)$ by $(\rho X,\rho Y)$ and let
$\rho\uparrow1$.  In \eqref{sp:doubleLi} this means damping its outer
color by $\rho$.  This convention is unambiguous for the combinations
appearing here.  Indeed
\[
 \mathcal T(0,B,C;X,Y)
 =\frac1{\Gamma(C)}\int_0^\infty
 t^{C-1}\operatorname{Li}_0(Xe^{-t})
                    \operatorname{Li}_B(Ye^{-t})\,dt
\]
for $\Re C>0$.  The product is analytic at $t=0$ and exponentially
decreasing at infinity.  Taylor subtraction at zero, followed by the
zeros of $1/\Gamma(C)$, gives an entire continuation in $(B,C)$.
The same argument applies with arbitrary fixed integer first orders,
and identifies the radial continuation locally uniformly in the
orders.  Thus all the order derivatives used below exist.

\begin{theorem}[Integer inner orders and all normal jets]
\label{sp:integer}
For positive integers $A,B$, one has the identity of entire functions
of $C$ at nontrivial unit colors
\begin{align}\label{sp:positive}
 \mathcal T(A,B,C;X,Y)
 ={}&\sum_{j=1}^{A}\binom{A+B-j-1}{B-1}
       \operatorname{Li}_{C+A+B-j,j}(Y,X/Y)\notag\\
 &+\sum_{j=1}^{B}\binom{A+B-j-1}{A-1}
       \operatorname{Li}_{C+A+B-j,j}(X,Y/X).
\end{align}
For $M\geq0$ an integer and arbitrary $B\in\mathbb C$,
\begin{equation}\label{sp:nonpositive}
 \mathcal T(-M,B,C;X,Y)
 =\sum_{j=0}^{M}(-1)^j\binom Mj
       \operatorname{Li}_{C-M+j,B-j}(X,Y/X).
\end{equation}
Every fixed $C$-derivative can be applied to these formulas.  In
\eqref{sp:nonpositive}, every fixed $B$-derivative is also valid.
Together with exchange of the two inner variables, these formulas
give a finite reduction for every pair of integer inner orders.
\end{theorem}
\begin{proof}
For positive $m,n$,
\[
 \frac1{m^A n^B}
 =\sum_{j=1}^{A}\binom{A+B-j-1}{B-1}
     \frac1{m^j(m+n)^{A+B-j}}
 +\sum_{j=1}^{B}\binom{A+B-j-1}{A-1}
     \frac1{n^j(m+n)^{A+B-j}}.
\]
It follows by repeated use of
$1/(mn)=(1/m+1/n)/(m+n)$, or by induction on $A+B$.
Multiply by the colors and $(m+n)^{-C}$ and put $k=m+n$.
The first sum has color $Y^k(X/Y)^m$, and the second has
$X^k(Y/X)^n$.  This proves \eqref{sp:positive} in an absolutely
convergent region.  For \eqref{sp:nonpositive}, expand
$(k-n)^M=\sum_j(-1)^j\binom Mj k^{M-j}n^j$ after the same reindexing.
Local normal convergence inside the bidisc and then the specified
analytic continuation prove both identities and the derivative claims.
\end{proof}

Here is the promised complete $A,B\in\{0,1\}$ table.  A superscript
$[r]$ on a double polylogarithm in this table differentiates only its
first order.  For $X\ne Y$,
\begin{align}
 \partial_C^r\mathcal T(0,0,C;X,Y)
 &=\frac{Y\operatorname{Li}_{C}^{[r]}(X)
          -X\operatorname{Li}_{C}^{[r]}(Y)}{X-Y},
 \label{sp:00}\\
 \partial_C^r\mathcal T(0,1,C;X,Y)
 &=\operatorname{Li}_{C,1}^{[r]}(X,Y/X),\label{sp:01}\\
 \partial_C^r\mathcal T(1,0,C;X,Y)
 &=\operatorname{Li}_{C,1}^{[r]}(Y,X/Y),\label{sp:10}\\
 \partial_C^r\mathcal T(1,1,C;X,Y)
 &=\operatorname{Li}_{C+1,1}^{[r]}(X,Y/X)
   +\operatorname{Li}_{C+1,1}^{[r]}(Y,X/Y).
 \label{sp:11}
\end{align}
The confluent version of \eqref{sp:00} is
\begin{equation}\label{sp:00confluent}
 \partial_C^r\mathcal T(0,0,C;z,z)
 =\operatorname{Li}_{C-1}^{[r]}(z)
    -\operatorname{Li}_{C}^{[r]}(z).
\end{equation}
All $C=-N$ and $C=0$ substitutions are now legitimate.  Formula
\eqref{sp:11} is a finite depth-two answer; it does not claim a
reduction of the remaining depth-two order derivatives to ordinary
polylogarithms.

\begin{remark}[The permitted derivative directions]
\label{sp:directionwarning}
In \eqref{sp:positive}, the first two orders are fixed integers.  Its
finite expression does not authorize differentiating with respect to
$A$ or $B$.  In particular, it does not reduce the full mixed jet
$\partial_A\partial_B\partial_C\mathcal T$ by formal differentiation
of the finite sums.  In \eqref{sp:nonpositive}, $B$ really is free,
so $B$-derivatives are allowed, while $A$ remains the fixed integer
$-M$.  This is the same distinction that prevents differentiating
a negative-$C$ value formula transversely without further information.
\end{remark}

\subsection{Rational generating functions give depth-one closure}

Let $d\geq1$, let $p_j$ be positive integers, let $a_j$ be nonnegative
integers, and let $X_j$ be nonzero complex numbers with $|X_j|\leq1$.
Define, initially for $\Re C>a_1+\cdots+a_d+d$,
\begin{equation}\label{sp:weighteddef}
 Z(C)=\sum_{m_1,\ldots,m_d\geq1}
  \frac{\prod_{j=1}^d m_j^{a_j}X_j^{m_j}}
       {(p_1m_1+\cdots+p_dm_d)^C}.
\end{equation}
Consider its coefficient generating function
\begin{equation}\label{sp:rationalR}
 R(t)=\prod_{j=1}^d\operatorname{Li}_{-a_j}(X_jt^{p_j})
      =\sum_{n\geq1}r_nt^n.
\end{equation}
The elementary identity
$\operatorname{Li}_{-a}(z)=(z\,d/dz)^a\{z/(1-z)\}$ proves that
$R$ is rational and bounded at infinity.  Write its partial fractions
in the form
\begin{equation}\label{sp:partialR}
 R(t)=c_\infty+
       \sum_{\alpha}\sum_{h=1}^{M_\alpha}
           \frac{c_{\alpha,h}}{(1-\alpha t)^h},
\end{equation}
where each $\alpha$ is a root of some equation
$\alpha^{p_j}=X_j$.  A permissible multiplicity bound is
\[
 M_\alpha=\sum_{j:\,\alpha^{p_j}=X_j}(a_j+1).
\]
This is an upper bound: zeros of other factors can cancel some poles.
In that event the corresponding highest coefficients are zero.
There is no positive-degree polynomial part because $R$ is bounded
at infinity.  An exact coefficient formula, also valid at repeated
roots, is
\begin{equation}\label{sp:cformula}
 c_{\alpha,h}
   =[u^{M_\alpha-h}]\,
       u^{M_\alpha}R((1-u)/\alpha).
\end{equation}
The function whose coefficient is extracted is holomorphic at $u=0$.

\begin{theorem}[Arbitrary depth, weights, multiplicities, and normal jets]
\label{sp:rationaltheorem}
Set
\begin{align}
 e_{h,r}
  &=[u^r]\prod_{\nu=1}^{h-1}(1+u/\nu),\label{sp:eh}\\
 d_{\alpha,r}
  &=\sum_{h=r+1}^{M_\alpha}c_{\alpha,h}e_{h,r}.
 \label{sp:dcoeff}
\end{align}
Empty products equal one.  Then
\begin{equation}\label{sp:depthone}
 Z(C)=\sum_\alpha\sum_{r=0}^{M_\alpha-1}
               d_{\alpha,r}\operatorname{Li}_{C-r}(\alpha).
\end{equation}
This is the meromorphic continuation of \eqref{sp:weighteddef}, and
every $C$-derivative is obtained by replacing each polylogarithm by
the corresponding order derivative.

If no $X_j$ equals one, the continuation is entire.  In general its
only possible poles are simple poles at $C=r+1$ from the $\alpha=1$
sector, with residue $d_{1,r}$.  For $k\geq0$, the constant Laurent
coefficient of $Z^{(k)}$ at an integer $m\geq1$ is
\begin{equation}\label{sp:stieltjesFP}
 \operatorname{CT}_{C=m} Z^{(k)}(C)
 =(-1)^k d_{1,m-1}\gamma_k
   +\sum_{(\alpha,r)\ne(1,m-1)}
      d_{\alpha,r}\operatorname{Li}_{m-r}^{[k]}(\alpha).
\end{equation}
Missing coefficients are zero; for $\alpha=1$ the nonsingular terms
mean ordinary Riemann-zeta derivatives.
\end{theorem}
\begin{proof}
The original multiple sum converges absolutely in the stated
half-plane.  Grouping it according to $n=\sum p_jm_j$ gives
$Z(C)=\sum_{n\geq1}r_n n^{-C}$.  The binomial series gives
\[
 [t^n](1-\alpha t)^{-h}
 =\alpha^n\binom{n+h-1}{h-1}
 =\alpha^n\prod_{\nu=1}^{h-1}(1+n/\nu)
 =\alpha^n\sum_{r=0}^{h-1}e_{h,r}n^r.
\]
For $n\geq1$ the constant $c_\infty$ contributes nothing.  Substitution
proves \eqref{sp:depthone} where the series converges absolutely and hence
by continuation.  If $\alpha\ne1$ and $|\alpha|\leq1$, the fixed-color
polylogarithm is entire in its order; for $\alpha=1$ it is $\zeta$.
The pole and derivative statements follow.  Finally
\[
 \zeta(1+v)=v^{-1}+
       \sum_{j\geq0}\frac{(-1)^j\gamma_j}{j!}v^j
\]
shows that the constant coefficient after $k$ differentiations is
$(-1)^k\gamma_k$.  The other terms are holomorphic at the specified
integer, proving \eqref{sp:stieltjesFP}.
\end{proof}

The harmonic-number coefficients in this theorem can be computed by
the finite-degree expansion
\begin{equation}\label{sp:harmonicexp}
 \prod_{\nu=1}^{h-1}(1+u/\nu)
 =\exp\!\left(
    \sum_{j\geq1}\frac{(-1)^{j+1}}j H_{h-1}^{(j)}u^j\right).
\end{equation}
For example $e_{h,0}=1$, $e_{h,1}=H_{h-1}$, and
$e_{h,2}=(H_{h-1}^2-H_{h-1}^{(2)})/2$.
All multiplicity corrections are consequently finite polynomials in
generalized harmonic numbers.  Coincidence in this statement concerns
the poles of the rational coefficient generating function.  It does
not identify two independently specified spatial finite parts.

\subsection{Roots of unity, Gamma values, and an explicit weighted example}

For a nontrivial $q$th root of unity $\alpha$, residue-class summation
gives
\begin{equation}\label{sp:DFT}
 \operatorname{Li}_s(\alpha)
   =q^{-s}\sum_{j=1}^q\alpha^j\zeta(s,j/q).
\end{equation}
It follows first where the series converge absolutely and then by
continuation.  In particular, for $s\ne1$,
\begin{equation}\label{sp:DFTjets}
 \operatorname{Li}_s^{[k]}(\alpha)
 =q^{-s}\sum_{j=1}^q\alpha^j
       \sum_{\ell=0}^k\binom{k}{\ell}(-\log q)^{k-\ell}
                 \zeta^{[\ell]}(s,j/q).
\end{equation}
At $s=1$, cancellation of the common pole gives instead
\begin{equation}\label{sp:DFTstieltjes}
 \operatorname{Li}_1^{[k]}(\alpha)
 =\frac1q\sum_{j=1}^q\alpha^j
       \sum_{\ell=0}^k\binom{k}{\ell}
        (-\log q)^{k-\ell}(-1)^\ell\gamma_\ell(j/q).
\end{equation}
The cancellation is performed before evaluation.  Applying Lerch's
formula $\zeta'(0,a)=\log\Gamma(a)-\tfrac12\log(2\pi)$
\cite[Eq.~25.11.18]{DLMFHzeta} gives
\begin{equation}\label{sp:Gammajet}
 \operatorname{Li}_0^{[1]}(\alpha)
   =\sum_{j=1}^q\alpha^j\log\Gamma(j/q)
       -\log q\,\frac{\alpha}{1-\alpha}.
\end{equation}
Thus simple poles in \eqref{sp:partialR} yield a finite log-Gamma
evaluation for $Z'(0)$ at roots of unity.  Repeated poles additionally
introduce the explicitly listed Hurwitz derivatives at negative
integers from \eqref{sp:DFTjets}; they should not be silently discarded.

For a concrete unequal-weight example, put
\[
 Z_{2,3}(C)=\sum_{m,n\geq1}\frac{(-1)^{m+n}}{(2m+3n)^C}
 \qquad(\Re C>2).
\]
Its rational generating function has the exact certificate
\begin{equation}\label{sp:weightedcertificate}
 \frac{t^5}{(1+t^2)(1+t^3)}
 =\frac{t^5-t^7-t^8+t^9+t^{10}-t^{12}}{1-t^{12}}.
\end{equation}
Consequently, for every $C$ by entire continuation,
\begin{align}\label{sp:weightedHurwitz}
 Z_{2,3}(C)=12^{-C}\bigl\{
  &\zeta(C,5/12)-\zeta(C,7/12)-\zeta(C,2/3)\notag\\
  &+\zeta(C,3/4)+\zeta(C,5/6)-\zeta(C,1)\bigr\}.
\end{align}
The six residues at $C=1$ cancel.  In particular,
\begin{equation}\label{sp:weightedGamma}
 Z_{2,3}'(0)
 =\log\frac{\Gamma(5/12)\Gamma(3/4)\Gamma(5/6)}
                 {\Gamma(7/12)\Gamma(2/3)}-\frac14\log12.
\end{equation}
An entirely convergent integral version, with no formal double sum
at $C=0$, is
\begin{align}\label{sp:weightedIntegral}
 &\int_0^1\left\{\frac1{(e^{2t}+1)(e^{3t}+1)}-\frac14\right\}
                         \frac{dt}{t}
 +\int_1^\infty\frac{dt}{t(e^{2t}+1)(e^{3t}+1)}
 +\frac\gamma4\notag\\
 &\hspace{24mm}=
 \log\frac{\Gamma(5/12)\Gamma(3/4)\Gamma(5/6)}
                 {\Gamma(7/12)\Gamma(2/3)}-\frac14\log12.
\end{align}
To verify the integral, insert the Mellin representation for
$\Re C>0$, subtract its constant $1/4$ at zero, and use
$1/\Gamma(C)=C+\gamma C^2+O(C^3)$.

For a fully confluent example, $d$ identical colors with unit weights
and zero inner orders give
\begin{equation}\label{sp:compositions}
 \sum_{m_1,\ldots,m_d\geq1}
       \frac{z^{m_1+\cdots+m_d}}{(m_1+\cdots+m_d)^C}
 =\frac1{(d-1)!}\sum_{r=0}^{d-1}s(d,r+1)
                         \operatorname{Li}_{C-r}(z),
\end{equation}
where $s(d,j)$ are the signed Stirling numbers of the first kind.
Indeed the coefficient at total index $n$ counts compositions and
equals $\binom{n-1}{d-1}$, including its zero values for $1\leq n<d$.
This proves the identity and all its normal derivatives.  For $d=3$,
$z=-1$, and differentiation at $C=0$, it yields
\begin{equation}\label{sp:triplealternating}
 \left.\frac{d}{dC}
 \sum_{m,n,k\geq1}\frac{(-1)^{m+n+k}}{(m+n+k)^C}
 \right|_{C=0}^{\mathrm{cont}}
 =-\frac{7\zeta(3)}{8\pi^2}
   -\frac92\zeta'(-1)-\frac12\log\pi.
\end{equation}
Here ``cont'' explicitly means continuation from $\Re C>3$.
The reduction uses
$\operatorname{Li}_s(-1)=(2^{1-s}-1)\zeta(s)$,
$\zeta'(0)=-\tfrac12\log(2\pi)$, and
$\zeta'(-2)=-\zeta(3)/(4\pi^2)$.

\subsection{An exact primitive hierarchy}

Keep the coefficients $r_n$ of \eqref{sp:rationalR} and put
$Z(C;z)=\sum_{n\geq1}r_n z^n/n^C$ for $|z|<1$.
The proof of Theorem~\ref{sp:rationaltheorem} gives
\[
 Z(C;z)=\sum_{\alpha,r}d_{\alpha,r}
                    \operatorname{Li}_{C-r}(\alpha z).
\]
For every nonnegative integer $k$,
\begin{equation}\label{sp:primitive}
 \int_0^z Z^{(k)}(C;t)\frac{dt}{t}=Z^{(k)}(C+1;z),
 \qquad
 z\frac{d}{dz}Z^{(k)}(C;z)=Z^{(k)}(C-1;z).
\end{equation}
The integrand is holomorphic at zero after division by $t$, and the
primitive is fixed by its zero value there.  These identities follow
termwise in the unit disc and continue along any common path avoiding
the polylogarithmic singularities.  This gives exact antiderivatives
of every normal jet without changing its spectral direction.

\subsection{Independent algebraic and numerical checks}

The accompanying script \texttt{verify\_spectral.py} performs 402
exact finite checks: 25 positive-order partial fractions, 12
composition polynomials, 105 harmonic--Stirling coefficients, seven
reconstructions of rational generating functions, 252 independently
convolved coefficients, and the cyclotomic certificate
\eqref{sp:weightedcertificate}.  The rational cases include unequal
weights, repeated roots, multiplicity six, and cancellation of an
otherwise apparent pole.  Nine independently evaluated convergent
Mellin integrals at 60 decimal digits check the first normal jets for
distinct colors $i,-i$, repeated color $-1$ in depths two and three,
the unequal weights $(2,3)$, and the root-Gamma formula for roots of
orders $3,4,5,6,12$.  The largest observed discrepancy is
$4.7\times10^{-61}$.  The exact algebraic checks and analytic proofs
are separate from these floating-point diagnostics; the latter are
not interval-certified error bounds.  Inputs, formulas, and full
outputs are supplied in \texttt{spectral\_results.json}.

% END sections/05_spectral.tex

% BEGIN sections/06_rays.tex
\section{The first two jets on every admissible ray through the origin}
\label{sp:sec:rays}

The uncolored Tornheim function is not jointly holomorphic at the
origin.  Consequently $\mathcal T(0,0,0;1,1)$ is not a value that can
be assigned independently of a restriction convention.  Borwein and
Dilcher established this explicitly for the family
$\mathcal T(t,t,\tau t;1,1)$ and evaluated its first derivative
\cite[Theorems 2.7 and 3.1]{BorweinDilcher}.  The diagonal value $1/3$ is also
recorded by Romik \cite[Theorem 1.3]{sp:Romik}, who credits the simple
diagonal derivative to Borwein and Dilcher.

We give the asymmetric first jet and an exact second-jet formula for
every ray avoiding the two polar directions.  Let
\[
 \mathcal T(A,B,C)=\mathcal T(A,B,C;1,1),\qquad L=\log(2\pi),
 \qquad W_{a,b,c}(t)=\mathcal T(at,bt,ct).
\]
Each restriction is formed away from the polar set and then continued
as a one-variable function.  It is not the evaluation of a nonexistent
joint Taylor series at $(0,0,0)$.

\begin{theorem}[Asymmetric ray values and quadratic jets]
\label{sp:raytheorem}
Let $a,b,c\in\mathbb C$ satisfy $(a+c)(b+c)\ne0$, and define
\begin{equation}\label{sp:rayB}
 \mathcal B(a,b,c)
 =\frac{c(c^2+ab-a^2-b^2)}{(a+c)(b+c)}.
\end{equation}
Then $W_{a,b,c}$ has a removable singularity at $t=0$, and
\begin{align}
 W_{a,b,c}(0)
   &=\frac14+\frac{c}{12(a+c)}+\frac{c}{12(b+c)},
       \label{sp:ray0}\\
 W_{a,b,c}'(0)
   &=\frac{a+b+2c}{4}L+\mathcal B(a,b,c)\zeta'(-1),
       \label{sp:ray1}\\
 W_{a,b,c}''(0)
   &=\frac{2ab+c(a+b)}4L^2
      -\frac{a^2+b^2+2c(a+b)+2c^2}{2}\zeta''(0)\notag\\
   &\hspace{13mm}+(a+b+c)\mathcal B(a,b,c)\zeta''(-1).
       \label{sp:ray2}
\end{align}
\end{theorem}

The equal-inner-order specialization is especially compact:
\begin{align}
 W_{1,1,\tau}(0)&=\frac{5\tau+3}{12(\tau+1)},\notag\\
 W_{1,1,\tau}'(0)&=\frac{\tau+1}{2}L
             +\frac{\tau(\tau-1)}{\tau+1}\zeta'(-1),\notag\\
 W_{1,1,\tau}''(0)&=\frac{\tau+1}{2}L^2
             -(\tau+1)^2\zeta''(0)
             +\frac{\tau(\tau-1)(\tau+2)}{\tau+1}\zeta''(-1).
 \label{sp:equalray}
\end{align}
The first two lines recover the cited classical result; the third is
the quadratic extension proved here.  In particular,
\begin{equation}\label{sp:diagonal2}
 \left.\frac{d^2}{dt^2}\mathcal T(t,t,t)\right|_{t=0}
    =\log^2(2\pi)-4\zeta''(0).
\end{equation}

\subsection{A local meromorphic decomposition}

Put $F(A,B;x)=\operatorname{Li}_A(e^{-x})
                         \operatorname{Li}_B(e^{-x})$.
Jonqui\`ere's convergent expansion for $0<x<2\pi$ gives
\[
 \operatorname{Li}_A(e^{-x})
  =\Gamma(1-A)x^{A-1}
    +\sum_{j\geq0}\zeta(A-j)\frac{(-x)^j}{j!}
\]
for $A$ near zero; all powers of $x$ use its real logarithm.
Subtract the following six terms:
\begin{align}\label{sp:localsix}
 S(A,B;x)={}&\Gamma(1-A)\Gamma(1-B)x^{A+B-2}
   +\Gamma(1-A)\zeta(B)x^{A-1}\notag\\
 &+\Gamma(1-B)\zeta(A)x^{B-1}
   -\Gamma(1-A)\zeta(B-1)x^A\notag\\
 &-\Gamma(1-B)\zeta(A-1)x^B+\zeta(A)\zeta(B).
\end{align}
On a sufficiently small parameter polydisc the remainder and each
fixed parameter derivative are $O(x^{1-\delta}(1+|\log x|^N))$ for
some $\delta<1$.  Consequently the function
\begin{align}\label{sp:Hdef}
 H(A,B,C)={}&\int_0^1 x^{C-1}\{F(A,B;x)-S(A,B;x)\}\,dx
            +\int_1^\infty x^{C-1}F(A,B;x)\,dx\notag\\
 &+\frac{\Gamma(1-A)\Gamma(1-B)}{A+B+C-2}
   +\frac{\Gamma(1-A)\zeta(B)}{A+C-1}
   +\frac{\Gamma(1-B)\zeta(A)}{B+C-1}
\end{align}
is jointly holomorphic near $(0,0,0)$ and symmetric in $A,B$.
The Mellin formula now gives the exact meromorphic decomposition
\begin{align}\label{sp:meromorphicH}
 \mathcal T(A,B,C)=\frac1{\Gamma(1+C)}\biggl\{
  &\zeta(A)\zeta(B)
   -\frac{C\Gamma(1-A)\zeta(B-1)}{A+C}\notag\\
  &-\frac{C\Gamma(1-B)\zeta(A-1)}{B+C}
   +C H(A,B,C)\biggr\}.
\end{align}
To justify it, integrate each term of \eqref{sp:localsix} explicitly
in a common Mellin convergence region, and then continue; the
subtracted integrals in \eqref{sp:Hdef} supply the required local
continuation.  The only polar hyperplanes meeting the origin are now
displayed.  On an admissible ray, their quotients are constants, so
the apparent one-variable singularity is removable.

\subsection{The diagonal input}

The diagonal identity \eqref{sp:diagonal2} has already been proved in
Theorem~\ref{cub:diagonal} by comparing the local cubic residue with
its Hurwitz Fourier representation. That comparison uses only the
coefficient of the simple pole: the unknown cubic finite part first
enters one order later. There is therefore no dependence on a numerical
value of the cubic integral in the argument below.

\subsection{Completing the asymmetric calculation}

The exact restriction
\begin{equation}\label{sp:axis}
 \mathcal T(0,0,C)=\zeta(C-1)-\zeta(C)
\end{equation}
follows by counting the $k-1$ positive pairs summing to $k$ in its
absolute-convergence half-plane.  Restricting
\eqref{sp:meromorphicH} to this axis gives
\[
 \Gamma(1+C)\{\zeta(C-1)-\zeta(C)\}
       =\frac5{12}+C H(0,0,C).
\]
Hence, writing all partial derivatives of $H$ below at the origin,
\begin{align}
 H(0,0,0)
   &=\zeta'(-1)+\frac L2-\frac{5\gamma}{12},\label{sp:H0}\\
 H_C
   &=\frac{\zeta''(-1)-\zeta''(0)}2
       -\gamma\left(\zeta'(-1)+\frac L2\right)
       +\frac5{24}\{\gamma^2+\zeta(2)\}.
       \label{sp:HC}
\end{align}
The symmetry $H_A=H_B$, together with \eqref{sp:diagonal2} in
\eqref{sp:meromorphicH}, then gives
\begin{equation}\label{sp:HA}
 H_A=H_B=\frac{L^2}{8}-\frac{\zeta''(0)}2
         +\gamma\zeta'(-1)-\frac{\gamma L}{4}
         -\frac{\gamma^2+\zeta(2)}{24}.
\end{equation}
These equations determine the entire constant and linear Taylor
polynomial of the jointly holomorphic remainder.

To prove Theorem~\ref{sp:raytheorem}, substitute
$A=at,B=bt,C=ct$ in \eqref{sp:meromorphicH}, use
\eqref{sp:H0}--\eqref{sp:HA}, and expand to second order.  The
constant is \eqref{sp:ray0}.  The coefficient of $\zeta'(-1)$ in
the first derivative is
\[
 c-\frac{bc}{a+c}-\frac{ac}{b+c}
       =\mathcal B(a,b,c),
\]
and all Euler-constant terms cancel, proving \eqref{sp:ray1}.
In the second derivative, the surviving coefficients of
$L^2,\zeta''(0),\zeta''(-1)$ are respectively
\[
 \frac{2ab+c(a+b)}4,\qquad
 -\frac{a^2+b^2+2c(a+b)+2c^2}{2},\qquad
 (a+b+c)\mathcal B(a,b,c).
\]
The coefficients of $\gamma^2$, $\zeta(2)$,
$\gamma\zeta'(-1)$ and $\gamma L$ cancel identically.  This proves
\eqref{sp:ray2} and completes the theorem.

For example, the diagonal and the third-order-parameter axis give
\[
 W_{1,1,1}(0)=\frac13,\qquad W_{0,0,1}(0)=\frac5{12}.
\]
For the asymmetric ray $(a,b,c)=(1,2,3)$ one obtains
\begin{align*}
 W_{1,2,3}(0)&=\frac{29}{80},\\
 W_{1,2,3}'(0)&=\frac94L+\frac9{10}\zeta'(-1),\\
 W_{1,2,3}''(0)&=\frac{13}{4}L^2-\frac{41}{2}\zeta''(0)
                             +\frac{27}{5}\zeta''(-1).
\end{align*}
These are explicit directional identities, and also a concrete audit
test for any routine that substitutes integer orders before
differentiating.  A routine returning the same value for the first
two restrictions has erased the specified limiting direction.

\paragraph{Independent numerical check.}
The accompanying \texttt{verify\_rays.py} integrates the convergent
Jonqui\`ere series on $(0,1)$ and an incomplete-Gamma double series on
$(1,\infty)$.  It does not evaluate a cubic finite-part formula.
At 65-digit working precision, seven successive symmetric step sizes
and polynomial extrapolation reproduce the first two derivatives on
the rays $(1,1,1)$, $(1,2,3)$, $(2,3,1)$, and $(1,-2,3)$.
The largest reported absolute discrepancy is below $2.6\times10^{-37}$.
These are diagnostics of the analytic proof, with truncation and
floating-point errors not certified by interval arithmetic.

% END sections/06_rays.tex

% BEGIN sections/07_harmonic.tex
% Self-contained section. Requires amsmath, amssymb, amsthm, hyperref.
% All labels/citation keys use h2: prefix. No custom macros required.
\section{Two independent exponents at a harmonic resonance}
\label{h2:sec:main}

The fully shifted height-one theorem in the incoming exact-identities
report varies the outer exponent and keeps every inner exponent equal
to one. Its final research section asks what survives when inner
exponents move independently. We answer a precise part of that question:
the first inner exponent may vary independently of the outer exponent,
with an arbitrary number of subsequent unit exponents. Every principal
Laurent coefficient and the finite part along every transverse affine
direction have a finite Gamma--Bernoulli formula. A polynomial in the
common shift supplies the same correction in all directions.

The meromorphic continuation of multiple zeta functions and the
dependence of specializations on the approach direction are classical
\cite{h2:AET,h2:MOW}. In particular, Matsumoto, Onozuka and Wakabayashi
give algorithms for Laurent expansions at arbitrary integer points,
allowing remainder integrals in general. The result below is a compact
closed coefficient formula for this particular common-shift family,
uniform in its trailing depth and its two-dimensional direction. We do
not claim priority for general continuation or multivariable Laurent
expansion. We also do not resolve arbitrary independent perturbations of
all trailing exponents.

\subsection{The family and its finite residue polynomial}

Throughout this section, $\Re a>0$, and powers of $n+a$ use the principal
logarithm. Define
\begin{align}
 e_r(n;a)&=[u^r]\prod_{j=0}^{n-1}\left(1+\frac{u}{j+a}\right),
 &p_n(a,u)&=\frac{(a+u)_n}{(a)_n},\label{h2:eq:e}\\
 Z_r(s,t;a)&=\sum_{n_0>n_1>\cdots>n_{r+1}\ge0}
 \frac{1}{(n_0+a)^s(n_1+a)^t\prod_{j=2}^{r+1}(n_j+a)},
 &&r\ge0.\label{h2:eq:Z}
\end{align}
The final product is empty at $r=0$. Thus $Z_0$ is a double Hurwitz
zeta function, and $Z_r(s,1;a)$ is the height-one function with $r+1$
inner unit exponents. Its depth generator is
\begin{equation}\label{h2:eq:generator}
 \mathcal Z_a(s,t,u)=\sum_{n\ge0}
     p_n(a,u)(n+a)^{-t}\zeta(s,n+a+1)
       =\sum_{r\ge0}Z_r(s,t;a)u^r.
\end{equation}
For example, absolute normal convergence holds when
$\Re s>1$ and $\Re(s+t)>2+\Re u$; a small closed disc in $u$
can be chosen in any strict version of these inequalities. The estimate
$p_n=O(n^{\Re u})$ and the usual Hurwitz-tail estimate prove this
claim and justify coefficient extraction. At each fixed depth,
$e_r(n;a)=O((1+\log n)^r)$ gives the corresponding normal convergence.

Let $B_k(x)$ denote Bernoulli polynomials, with $B_1(0)=-1/2$.
Define polynomials $q_j(u)$ through the formal series
\begin{equation}\label{h2:eq:qgen}
 \sum_{j\ge0}q_j(u)z^j
 =\exp\left(\sum_{k\ge1}
       \frac{(-1)^{k+1}\{B_{k+1}(u)-B_{k+1}(0)\}}
            {k(k+1)}z^k\right).
\end{equation}
This is a coefficient definition; no convergence of the infinite
asymptotic series is asserted. The classical Gamma-ratio expansion
\cite{DLMFGamma} reads
\begin{equation}\label{h2:eq:ratio}
 \frac{\Gamma(x+u)}{\Gamma(x)}
   \sim x^u\sum_{j\ge0}q_j(u)x^{-j}.
\end{equation}
In particular,
\[
 \begin{aligned}
 q_0&=1,& q_1&=\frac{u(u-1)}2,\\
 q_2&=\frac{u(u-1)(u-2)(3u-1)}{24},&
 q_3&=\frac{u^2(u-1)^2(u-2)(u-3)}{48}.
 \end{aligned}
\]
Every finite truncation has a locally uniform remainder in $u$.
Put
\begin{align}
 C_a(u)&=\frac{\Gamma(a)}{\Gamma(a+u)},\label{h2:eq:C}\\
 A_0(s)&=\frac1{s-1},\quad A_1(s)=-\frac12,\quad
 A_{2k}(s)=\frac{B_{2k}}{(2k)!}(s)_{2k-1}\ (k\ge1),\quad
 A_{2k+1}(s)=0\ (k\ge1),\label{h2:eq:A}\\
 P_m(s,u)&=\sum_{j=0}^{m+1}q_j(u)A_{m+1-j}(s).
 \label{h2:eq:P}
\end{align}
Thus $P_m$ is a finite rational expression, analytic for $s$ near $-m$.
Equivalently,
\[
 P_m(s,u)=\frac{q_{m+1}(u)}{s-1}-\frac{q_m(u)}2
  +\sum_{k=1}^{\lfloor(m+1)/2\rfloor}
       \frac{B_{2k}(s)_{2k-1}}{(2k)!}q_{m+1-2k}(u).
\]
The first two examples are
\begin{align}
 P_0(s,u)&=\frac{u(u-1)}{2(s-1)}-\frac12,
 \label{h2:eq:P0}\\
 P_1(s,u)&=\frac{u(u-1)(u-2)(3u-1)}{24(s-1)}
             -\frac{u(u-1)}4+\frac{s}{12}.
 \label{h2:eq:P1}
\end{align}

\begin{proposition}[One moving divisor before depth extraction]
\label{h2:prop:local}
For every integer $m\ge0$, there is a jointly holomorphic function
$\mathcal H_{m,a}$ near $(s,t,u)=(-m,1,0)$ such that
\begin{equation}\label{h2:eq:local}
 \mathcal Z_a(s,t,u)=
 \frac{C_a(u)P_m(s,u)}{s+t+m-1-u}+\mathcal H_{m,a}(s,t,u).
\end{equation}
The statement is also locally uniform in $a$ on the right half-plane.
\end{proposition}
\begin{proof}
Write $x=n+a$. Since
$\zeta(s,x+1)=\zeta(s,x)-x^{-s}$, the Hurwitz asymptotic expansion
\cite{DLMFHzeta} gives
\[
 \zeta(s,x+1)\sim\sum_{k\ge0}A_k(s)x^{1-s-k}.
\]
Multiplying a finite truncation by \eqref{h2:eq:ratio} yields
\[
 p_n(a,u)x^{-t}\zeta(s,x+1)
 \sim C_a(u)\sum_{\ell\ge0}
   \left(\sum_{j+k=\ell}q_j(u)A_k(s)\right)
                         x^{1-s-t+u-\ell}.
\]
Subtract sufficiently many terms from the summand in
\eqref{h2:eq:generator}. The remainder series converges normally in a
neighborhood of $(-m,1,0)$: after retaining $L$ terms, its
summands satisfy
\[
 O(n^{1-\Re(s+t)+\Re u-L}).
\]
Choosing $L>m+2$ suffices on a small neighborhood. Parameter derivatives introduce only fixed powers
of $\log n$ after shrinking that neighborhood.

The sums of the subtracted terms are the corresponding finite linear
combination of
$\zeta(s+t-1-u+\ell,a)$. Exactly one can have a pole in this
neighborhood: $\ell=m+1$. Its residue is
$C_a(u)P_m(s,u)$, and subtracting its simple pole leaves a holomorphic
function. This proves \eqref{h2:eq:local}. The use of finite asymptotic
subtraction, with a normally convergent remainder, is what makes this a
meromorphic identity rather than a formal asymptotic assertion.
\end{proof}

\subsection{An explicit value of the holomorphic remainder}

The next lemma is the specialization mechanism from the incoming
fully shifted height-one theorem \cite{IncomingExact}. We include its
proof because the direction-independent correction in the main theorem
depends on it.

\begin{lemma}[A finite Stirling specialization]
\label{h2:lem:stirling}
Let
\[
 F_a(w,u)=\sum_{n\ge0}\frac{p_n(a,u)}{(n+a)^w},\qquad
 T_m(a,u)=-\sum_{j=1}^{m+1}
   \left\{\begin{matrix}m+1\\j\end{matrix}\right\}
      \frac{(a-1)^{\underline j}}{u+j}.
\]
Here $(a-1)^{\underline j}$ is a falling factorial and braces denote
Stirling numbers of the second kind. For sufficiently small nonzero
$u$, the meromorphic continuation satisfies
\begin{equation}\label{h2:eq:Fnegative}
 F_a(-m,u)=T_m(a,u).
\end{equation}
Consequently, with
\begin{equation}\label{h2:eq:L}
 L_m(a,u)=\frac{T_m(a,u)-\zeta(-m,a)}u,
\end{equation}
one has
\begin{equation}\label{h2:eq:Hvalue}
 \mathcal H_{m,a}(-m,1,u)
   =L_m(a,u)+\frac{C_a(u)P_m(-m,u)}u.
\end{equation}
The right side of \eqref{h2:eq:Hvalue} is holomorphic at $u=0$.
\end{lemma}
\begin{proof}
The Mellin kernel of $F_a$ is
\[
 G_a(v,u)=e^{-av}\,{}_2F_1(1,a+u;a;e^{-v}).
\]
Gauss connection at one gives
\begin{align*}
 G_a(v,u)&=\frac{\Gamma(a)\Gamma(1+u)}{\Gamma(a+u)}
       \frac{e^{-v}}{(1-e^{-v})^{1+u}}+R_a(v,u),\\
 R_a(v,u)&=-\frac{a-1}{1+u}e^{-av}
      {}_2F_1(1,a+u;2+u;1-e^{-v}).
\end{align*}
The second kernel is analytic at $v=0$ and decays exponentially at
infinity, locally in the parameters. Mellin subtraction therefore
shows that its quotient by $\Gamma(w)$ is entire in $w$, and has value
$(-1)^m m![v^m]R_a(v,u)$ at $w=-m$. The first kernel has local
exponents $-1-u+j$; for small nonzero $u$ its Mellin transform has no
pole at $w=-m$, so division by $\Gamma(w)$ makes its value zero there.
This proves that the left side of \eqref{h2:eq:Fnegative} is a
polynomial in $a$ of degree at most $m+1$.

To identify the polynomial, take a positive integer $a=A$. Direct
index shifting in the original Dirichlet series gives
\[
 F_A(w,u)=\frac{\Gamma(A)\Gamma(1+u)}{\Gamma(A+u)}
 \left(F_1(w,u)-\sum_{k=1}^{A-1}
       \frac{(1+u)_{k-1}}{(k-1)!k^w}\right).
\]
At $w=-m$, the first term vanishes. The finite identity
\[
 k^m=\sum_{j=1}^{m+1}
    \left\{\begin{matrix}m+1\\j\end{matrix}\right\}
                           (k-1)^{\underline{j-1}}
\]
and the binomial summation
\[
 \frac{\Gamma(A)\Gamma(1+u)}{\Gamma(A+u)}
 \sum_{k=1}^{A-1}\frac{(1+u)_{k-1}}{(k-1)!}
           (k-1)^{\underline{j-1}}
       =\frac{(A-1)^{\underline j}}{u+j}
\]
give $T_m(A,u)$. Both sides are polynomials in $a$ of degree at most
$m+1$, and agreement at all positive integers proves
\eqref{h2:eq:Fnegative} for every $\Re a>0$.

Finally, $p_{n+1}-p_n=up_n/(n+a)$ implies, first by absolutely
convergent summation,
\begin{equation}\label{h2:eq:telescoping}
 \mathcal Z_a(s,1,u)=\frac{F_a(s,u)-\zeta(s,a)}u.
\end{equation}
Meromorphic continuation, \eqref{h2:eq:local}, and
\eqref{h2:eq:Fnegative} now give \eqref{h2:eq:Hvalue}. Its
holomorphy also follows directly from Proposition~\ref{h2:prop:local}.
\end{proof}

The order of operations in \eqref{h2:eq:Fnegative} matters. It restricts
$w$ to $-m$ for nonzero $u$ and then continues the resulting function in
$u$. At a crossing of a moving pole it does not give the Laurent finite
part obtained by first extracting a coefficient in $u$. The main
theorem keeps these operations in the required order.

For later use put
\begin{equation}\label{h2:eq:Delta}
 \Delta_{m,d}(a)=(-1)^{d+1}\sum_{j=1}^{m+1}
  \left\{\begin{matrix}m+1\\j\end{matrix}\right\}
                      \frac{(a-1)^{\underline j}}{j^{d+1}}.
\end{equation}
Thus $[u^r]L_m(a,u)=\Delta_{m,r+1}(a)$ for $r\ge0$.

\subsection{The complete two-direction finite-part theorem}

\begin{theorem}[All depths and every transverse affine direction]
\label{h2:thm:direction}
Let $m,r\ge0$, $\Re a>0$, and $\alpha,\beta\in\mathbb C$ with
$q=\alpha+\beta\ne0$. Set $\theta=\alpha/q$ and define the analytic
germ
\begin{equation}\label{h2:eq:V}
 V_{m,\theta}(u;a)=C_a(u)P_m(-m+\theta u,u).
\end{equation}
Then
\begin{equation}\label{h2:eq:main}
 \boxed{\displaystyle
 Z_r(-m+\alpha\varepsilon,1+\beta\varepsilon;a)
 =\sum_{h=0}^{r+1}
     \frac{[u^{r+1-h}]V_{m,\theta}(u;a)}
          {q^h\varepsilon^h}
      +\Delta_{m,r+1}(a)+O(\varepsilon).}
\end{equation}
In particular,
\begin{equation}\label{h2:eq:FP}
 \operatorname{FP}_{\varepsilon=0}
 Z_r(-m+\alpha\varepsilon,1+\beta\varepsilon;a)
 =[u^{r+1}]V_{m,\theta}(u;a)+\Delta_{m,r+1}(a).
\end{equation}
Every displayed coefficient is a finite expression in rational numbers,
$a$, $\theta$, $q^{-1}$, and the polygamma values
\[
 \psi(a),\psi'(a),\ldots,\psi^{(r)}(a).
\]
\end{theorem}
\begin{proof}
In \eqref{h2:eq:local}, substitute the affine path and put
$f(\varepsilon,u)=C_a(u)P_m(-m+\alpha\varepsilon,u)$. The singular
term becomes $f(\varepsilon,u)/(q\varepsilon-u)$. For fixed nonzero
$\varepsilon$, expand in $u$ first:
\[
 \frac1{q\varepsilon-u}
    =\sum_{n\ge0}\frac{u^n}{q^{n+1}\varepsilon^{n+1}}.
\]
If $h\ge1$, collecting the coefficient of
$u^r\varepsilon^{-h}$ gives
\[
 q^{-h}[u^{r+1-h}]
       C_a(u)P_m\left(-m+\frac{\alpha}{q}u,u\right).
\]
For the constant Laurent coefficient the same calculation excludes the
term of degree zero in $\varepsilon$ in $f$, and therefore gives
\[
 [u^{r+1}]C_a(u)
   \left(P_m(-m+\theta u,u)-P_m(-m,u)\right).
\]
The holomorphic remainder contributes
$[u^r]\mathcal H_{m,a}(-m,1,u)$ to the constant coefficient and
nothing to the principal part. Formula \eqref{h2:eq:Hvalue} cancels
$-[u^{r+1}]C_a(u)P_m(-m,u)$ and leaves
$[u^r]L_m=\Delta_{m,r+1}$. This proves \eqref{h2:eq:main}.

Finally,
\[
 C_a(u)=\exp\left(-\sum_{k\ge1}
                  \frac{\psi^{(k-1)}(a)}{k!}u^k\right).
\]
Only finitely many coefficients, through degree $r+1$, enter the
formula. This proves the stated finite algebra and completes the proof.
\end{proof}

The choice $(\alpha,\beta)=(1,0)$ recovers the archived common-shift
height-one expansion. Indeed comparison of the moving residues in
\eqref{h2:eq:telescoping} proves the finite identity
\begin{equation}\label{h2:eq:recovery}
 uP_m(-m+u,u)=q_{m+1}(u).
\end{equation}
If the former notation is used, $q_{m+1}$ is precisely its polynomial
$W_m$. The choice $(\alpha,\beta)=(0,1)$ instead keeps the outer
exponent at $-m$ and moves the first inner exponent. Formula
\eqref{h2:eq:FP} specifies their discrepancy at every depth; it is
not legitimate to transfer a finite part from one direction to the
other without that correction.

\begin{corollary}[Pole order is independent of the transverse direction]
\label{h2:cor:poles}
For $m=0$ or positive odd $m$, the pole in \eqref{h2:eq:main} has order
exactly $r+1$, with leading coefficient
$\zeta(-m)/q^{r+1}$. For positive even $m$, it has order exactly $r$
when $r\ge1$, with leading coefficient
\[
 \frac{m+1}{2m}\frac{B_m}{q^r}.
\]
At positive even $m$ and $r=0$ the function is holomorphic at the
intersection, and the value is independent of the direction.
\end{corollary}
\begin{proof}
From \eqref{h2:eq:qgen}, $q_0(0)=1$ and $q_j(0)=0$ for $j>0$.
Hence $P_m(-m,0)=A_{m+1}(-m)=\zeta(-m)$. This proves the first
assertion. For positive even $m$, $A_{m+1}$ is identically zero, so
$[u]P_m(-m+\theta u,u)$ is independent of $\theta$. It can be
computed at $\theta=1$ using \eqref{h2:eq:recovery}.

For completeness, the Gamma-ratio coefficients also obey
\[
 q_{m+1}(u)=(u-m)_{m+1}[v^{m+1}]
                e^{-v}\left(\frac{v}{1-e^{-v}}\right)^{1+u}.
\]
To justify this auxiliary identity, the beta integral gives, for
$\Re u<0$ and sufficiently large positive $x$,
\[
 \frac{\Gamma(x+u)}{\Gamma(x)}
 =\frac1{\Gamma(-u)}\int_0^\infty e^{-xv}v^{-1-u}
       e^{-uv}\left(\frac{v}{1-e^{-v}}\right)^{1+u}\,dv.
\]
Expanding the analytic factor at $v=0$ and integrating a finite Taylor
polynomial gives
\[
 q_j(u)=(-u)_j[v^j]\,
          e^{-uv}\left(\frac{v}{1-e^{-v}}\right)^{1+u}.
\]
The analytic factor here is obtained from
$e^{-v}(v/(1-e^{-v}))^{1+u}$ by replacing $v$ with $-v$.
Since $(-1)^j(-u)_j=(u-j+1)_j$, this proves the identity with
$j=m+1$. Both sides are polynomials in $u$, so it holds for every $u$
by polynomial continuation.
For even $m\ge2$, the coefficient at $u=0$ in the last bracket
vanishes, while its coefficient of $u$ is
\[
 [v^{m+1}]\frac{v}{e^v-1}
                   \log\frac{v}{1-e^{-v}}
       =\frac{m+1}{2m}\frac{B_m}{m!}.
\]
This follows on multiplying the ordinary Bernoulli series by
$v/2-\sum_{k\ge1}B_{2k}v^{2k}/(2k(2k)!)$.
Since $(u-m)_{m+1}=m!u+O(u^2)$, one obtains
$[u^2]q_{m+1}=(m+1)B_m/(2m)\ne0$. Formula
\eqref{h2:eq:main} now gives the stated order and leading coefficient.
For $r=0$, the local numerator $P_m(s,0)=A_{m+1}(s)$ vanishes
identically, proving joint holomorphy.
\end{proof}

\paragraph{Directions within the polar hyperplane.}
The condition $q\ne0$ is essential whenever the pole in
Corollary~\ref{h2:cor:poles} is genuine. A nonconstant affine path with
$\alpha+\beta=0$ remains inside $s+t+m-1=0$, so the meromorphic function
has no ordinary one-variable restriction along that path. The exception
is $r=0$ and positive even $m$, where the function is jointly holomorphic:
all approach directions, including tangent ones, give
\[
 Z_0(-m,1;a)=\Delta_{m,1}(a)+\frac{m+1}{2m}B_m.
\]

\subsection{Explicit identities and direction corrections}

Write $P=\psi(a)$ and $Q=\psi'(a)$. Formula \eqref{h2:eq:P0}
gives
\[
 P_0(\theta u,u)=-\frac12+\frac u2
       +\frac{\theta-1}{2}u^2
       +\frac{\theta^2-\theta}{2}u^3+O(u^4).
\]
Consequently,
\begin{align}
 Z_0(\alpha\varepsilon,1+\beta\varepsilon;a)
 &= -\frac1{2q\varepsilon}+a-\frac12+\frac P2
                   +O(\varepsilon),\label{h2:eq:example00}\\
 Z_1(\alpha\varepsilon,1+\beta\varepsilon;a)
 &= -\frac1{2q^2\varepsilon^2}
      +\frac{1+P}{2q\varepsilon}
      +1-a+\frac{\theta-1}{2}-\frac P2
             -\frac{P^2-Q}{4}+O(\varepsilon).
 \label{h2:eq:example01}
\end{align}
Thus at total depth three the finite part changes by exactly
$(\theta-\eta)/2$ when the direction parameter changes from $\eta$
to $\theta$; the result holds at every common shift $a$.

At the next outer integer one obtains
\begin{equation}\label{h2:eq:example10}
 Z_0(-1+\alpha\varepsilon,1+\beta\varepsilon;a)
 =-\frac1{12q\varepsilon}
    +\frac{a^2+a}{4}+\frac{2\theta-5}{24}+\frac P{12}
                  +O(\varepsilon).
\end{equation}
The two coordinate directions differ by $1/12$ in their finite parts.
For an example where the highest potential pole cancels, put
$L=\gamma+2\log2$. At the half shift,
\begin{equation}\label{h2:eq:example21}
 Z_1(-2+\alpha\varepsilon,1+\beta\varepsilon;1/2)
    =\frac1{8q\varepsilon}+\frac L8-\frac{7+12\theta}{288}
                           +O(\varepsilon).
\end{equation}
At $\theta=1$ this is the previously obtained height-one finite part
$L/8-19/288$; at $\theta=0$ it is $L/8-7/288$.

In general define the difference between two directional finite parts
by $\mathfrak A_{m,r}(\theta,\eta;a)$. It has the exact form
\begin{equation}\label{h2:eq:anomaly}
 \mathfrak A_{m,r}(\theta,\eta;a)
 =[u^{r+1}]C_a(u)
       \bigl(P_m(-m+\theta u,u)-P_m(-m+\eta u,u)\bigr).
\end{equation}
The Stirling polynomial has canceled completely. At double-zeta depth
$r=0$ this reduces to
\begin{equation}\label{h2:eq:anomaly0}
 \mathfrak A_{m,0}(\theta,\eta;a)
   =-(\theta-\eta)\zeta(-m)H_m,
 \qquad H_0=0.
\end{equation}
To see the last equality, $[u]$ of the difference is
$(\theta-\eta)A'_{m+1}(-m)$. For odd $m$,
$A_{m+1}(s)=B_{m+1}(s)_m/(m+1)!$, and its logarithmic derivative
at $s=-m$ is $-H_m$; the cases $m=0$ and positive even $m$ vanish.
Thus the double-zeta directional correction is rational and independent
of $a$, even though each individual finite part contains $\psi(a)$.

\subsection{An exact normal slice and all Stieltjes derivatives}

At depth two one can go beyond the finite part in the direction that
keeps the outer exponent fixed. The resulting reduction is a classical
Bernoulli specialization of a multiple zeta function; it also supplies
particularly short mixed Gamma--Stieltjes identities.

\begin{theorem}[The complete inner-exponent slice]
\label{h2:thm:slice}
For $m\ge0$, the meromorphic identity
\begin{equation}\label{h2:eq:slice}
 Z_0(-m,t;a)=-\frac1{m+1}\sum_{j=0}^{m+1}
       \binom{m+1}{j}B_j(1)\zeta(t-m-1+j,a)
\end{equation}
holds. More generally, for every trailing depth $r$,
\begin{equation}\label{h2:eq:slicedepth}
 Z_r(-m,t;a)=-\frac1{m+1}\sum_{j=0}^{m+1}
       \binom{m+1}{j}B_j(1)M_r(t-m-1+j;a),
\end{equation}
where $M_r(w;a)=\sum_{n\ge0}e_r(n;a)(n+a)^{-w}$.
\end{theorem}
\begin{proof}
In the half-plane $\Re t>m+2$, the sum
$\sum e_r(n;a)(n+a)^{-t}\zeta(s,n+a+1)$ is locally normally
convergent for $s$ near $-m$. The Bernoulli evaluation
\[
 \zeta(-m,x+1)=-\frac{B_{m+1}(x+1)}{m+1}
 =-\frac1{m+1}\sum_{j=0}^{m+1}
       \binom{m+1}{j}B_j(1)x^{m+1-j}
\]
gives the finite formula termwise there. Continuing meromorphically in
$t$ proves the identities. This argument specifies the restriction
$s=-m$ before taking a Laurent coefficient in $t$.
\end{proof}

Use the Stieltjes convention
\[
 \zeta(1+\varepsilon,a)=\frac1\varepsilon
       +\sum_{\ell\ge0}\frac{(-1)^\ell}{\ell!}
                         \gamma_\ell(a)\varepsilon^\ell,
 \qquad \gamma_0(a)=-\psi(a).
\]
Writing $\zeta^{(\ell)}(w,a)=\partial_w^\ell\zeta(w,a)$,
\eqref{h2:eq:slice} yields the complete regular coefficients
\begin{align}
 [\varepsilon^\ell]Z_0(-m,1+\varepsilon;a)
 ={}&-\frac1{(m+1)\ell!}\sum_{j=0}^{m}
       \binom{m+1}{j}B_j(1)\zeta^{(\ell)}(j-m,a)
       \notag\\
 &+\frac{(-1)^\ell}{\ell!}\zeta(-m)\gamma_\ell(a),
 \qquad \ell\ge0.\label{h2:eq:slicejets}
\end{align}
Its principal part is $\zeta(-m)/\varepsilon$. At $m=0$ this simplifies to
\begin{equation}\label{h2:eq:allstieltjes}
 \boxed{\displaystyle
 \operatorname{FP}_{t=1}\partial_t^\ell Z_0(0,t;a)
   =-\zeta^{(\ell)}(0,a)
           +\frac{(-1)^{\ell+1}}2\gamma_\ell(a),
 \qquad \ell\ge0.}
\end{equation}
For $\ell=1$, Lerch's classical identity
\cite{DLMFHzeta} gives the concrete logarithmic formula
\begin{equation}\label{h2:eq:firststieltjes}
 \boxed{\displaystyle
 \operatorname{FP}_{t=1}\partial_t Z_0(0,t;a)
   =-\log\Gamma(a)+\frac12\log(2\pi)+\frac12\gamma_1(a).}
\end{equation}
The logarithm of $\Gamma$ is its analytic branch on $\Re a>0$,
real for positive $a$. Equivalently,
\begin{align}
 Z_0(0,1+\varepsilon;a)
  ={}&-\frac1{2\varepsilon}+a-\frac12+\frac12\psi(a)\notag\\
  &+\varepsilon\left(-\log\Gamma(a)+\frac12\log(2\pi)
                                     +\frac12\gamma_1(a)\right)
       +O(\varepsilon^2).\label{h2:eq:Z0next}
\end{align}
This does not evaluate the first regular coefficient in the different
slice $Z_0(\varepsilon,1;a)$; the two slices coincide only at those
coefficients explicitly identified above.

There is also a finite antiderivative at every Stieltjes order. Put
$J_\ell(a)=\operatorname{FP}_{t=1}\partial_t^\ell Z_0(0,t;a)$.
Then
\begin{equation}\label{h2:eq:primitive}
 \boxed{\displaystyle
 \int J_\ell(a)\,da
  =\frac{\zeta^{(\ell+1)}(0,a)}{2(\ell+1)}
      -\ell!\sum_{j=0}^{\ell}
                  \frac{\zeta^{(j)}(-1,a)}{j!}
       +\text{constant}.}
\end{equation}
Indeed the Hurwitz shift derivative implies
\[
 \partial_a\zeta^{(j)}(-1,a)
     =\zeta^{(j)}(0,a)-j\zeta^{(j-1)}(0,a),\qquad
 \partial_a\zeta^{(\ell+1)}(0,a)
     =(-1)^{\ell+1}(\ell+1)\gamma_\ell(a).
\]
The finite sum telescopes on differentiation. For example, a primitive
of the right side of \eqref{h2:eq:firststieltjes} is
\[
 \frac14\zeta''(0,a)-\zeta'(-1,a)-\zeta(-1,a).
\]
These formulas provide exact derivative and antiderivative identities
in the independent inner spectral direction, including every
generalized Stieltjes order.

\subsection{Scope, verification, and the next questions}

The new finite-part theorem stops at $O(\varepsilon)$ because the
first derivatives of $\mathcal H_{m,a}$ can enter that term.
The normal slice at depth two is exceptional: the Bernoulli reduction
evaluates all of its regular coefficients. Formula
\eqref{h2:eq:slicedepth} reduces the corresponding problem at larger
depth by one, but does not by itself reduce every resulting spectral
derivative to ordinary Hurwitz data.

The supplied program \texttt{verify\_two\_exponent.py} checks the
residue identity \eqref{h2:eq:recovery} exactly for $0\le m\le6$.
It separately evaluates the continued nested sum using its original
finite initial segment and a long Euler--Maclaurin tail. Discrete Cauchy
averages then extract Laurent coefficients without using the predicted
principal part in the extraction. A second cutoff and longer tail
check numerical stability. Eight cases at $65$ decimal working digits
compare $21$ Laurent coefficients through the finite part; the maximum
absolute discrepancy is below $2.25\times10^{-46}$. The longer-tail
repeat differs by less than $5.9\times10^{-52}$. These are
arbitrary-precision diagnostics;
the normal-convergence and coefficient-extraction proofs above establish
the identities for all parameters.

Three further research directions are concrete. First, allow the
second inner exponent to move as well; this introduces several
intersecting polar divisors and tests whether a finite analogue of
$P_m$ still eliminates every constant-term remainder. Second, determine
linear combinations of first regular coefficients in different
transverse directions that cancel the derivatives of
$\mathcal H_{m,a}$. Third, compare the rational directional correction
at $a=1$ with the general Laurent algorithms and with the
Stirling-polynomial formulas for nonpositive multiple-zeta values
\cite{h2:MOW,h2:IshiiShinohara}; this requires matching both the order of
indices and the chosen regularization, since the present base point has
unit inner exponents.

% END sections/07_harmonic.tex

% BEGIN sections/08_audit.tex
\section{Source audit and proposed corrections}
\label{audit:sec}

\subsection{The Gaussian candidates remain conjectures}

The manuscript distinguishes proved reductions from frozen
integer-relation candidates. That distinction is preserved here.
The current $S_6$ statement is
\texttt{cycloquot:conj:S6} in
\texttt{chapters/04-cyclotomic-quotients.tex}; the current $S_8$
statement is \texttt{s8new:conj:S8} in
\texttt{chapters/04-S8-candidate.tex} \cite{Repo}.
The latter is a different vector from an earlier rejected $S_8$
candidate. An objection to that rejected vector is not an objection
to the current one.

The supplied rigorous proximity certificates for the current frozen
vectors have residual intervals containing zero. In particular, a
bound $|R|<10^{-355}$ proves a very strong approximation but does not
prove $R=0$. Repeating floating-point integer-relation searches would
not close that logical gap. Nothing proved in this article supplies
the missing exact cyclotomic relation, so neither conjecture is
promoted to a theorem.

The finite separating functional already constructed for a restricted
level-four quotient is useful but has a limited conclusion. It rules
out a derivation inside the particular stated depth-one-product
shuffle/stuffle presentation. It does not prove that the target
constant lies outside every larger cyclotomic algebra, and it does
not rule out octahedral functional equations or additional
regularized relations. The exact proof of $S_4$ in the manuscript,
with its explicitly verified additional relation families, is a
model for the type of enlargement that would be needed.

% BEGIN sections/08a_s6_bridge.tex
\subsection{A convergent depth-three bridge for the Gaussian $S_6$ target}
\label{audit:s6-bridge}

The restricted quotient can be tested against an explicit relation just
outside its defining schemas. This gives a concrete exact reduction to
include in a future certificate, while leaving the $S_6$ conjecture open.
Here the manuscript's harmonic sum and Dirichlet beta function are
\[
 S_p=\sum_{n\ge0}\frac{(-1)^nH_n}{(2n+1)^p}\quad(p>1),
 \qquad
 \beta(s)=\sum_{n\ge0}\frac{(-1)^n}{(2n+1)^s}.
\]
Use the series convention
\[
 \Li_{s_1,\ldots,s_d}(x_1,\ldots,x_d)
 =\sum_{n_1>\cdots>n_d\ge1}
   \frac{x_1^{n_1}\cdots x_d^{n_d}}
        {n_1^{s_1}\cdots n_d^{s_d}},
 \qquad K_{6,1}=\operatorname{Im}\Li_{6,1}(i,-1).
\]
For precision about colors, write $x_j=i^{c_j}$, with $c_j$ read modulo
four. With $\omega_a=dt/(t-a)$ and outer-letter-first iterated integrals
$\mathcal I$, the corresponding word convention is
\begin{equation}
 \Li_{\boldsymbol s}(i^{c_1},\ldots,i^{c_d})
 =(-1)^d\mathcal I\left(
 0^{s_1-1}i^{-c_1}\cdots
 0^{s_d-1}i^{-(c_1+\cdots+c_d)}\right).
 \label{audit:s6-word-convention}
\end{equation}
All words below have first letter different from $1$ and last letter
different from $0$.

\begin{proposition}[A single-by-double shuffle reduction]
\label{audit:s6-depth-three}
The following identity consists entirely of convergent values:
\begin{align}
 K_{6,1}+\operatorname{Im}\Li_{5,2}(1,-i)
 ={}&\sum_{a=1}^{5}
       \operatorname{Im}\Li_{a,6-a,1}(i,-i,-1)\notag\\
 &+\operatorname{Im}\Li_{5,1,1}(1,i,i)
   +\operatorname{Im}\Li_{5,1,1}(1,-1,-i)\notag\\
 &-\operatorname{Im}\Li_{1,5,1}(i,1,-1)
   -\operatorname{Im}\Li_{5,1,1}(1,i,-1)\notag\\
 &-\operatorname{Im}\Li_{5,1,1}(1,-1,i).
 \label{audit:s6-depth-three-identity}
\end{align}
\end{proposition}
\begin{proof}
Consider $P=\Li_1(i)\Li_{5,1}(1,-1)$. Its stuffle expansion is
\begin{align*}
 P={}&\Li_{1,5,1}(i,1,-1)+\Li_{5,1,1}(1,i,-1)
       +\Li_{5,1,1}(1,-1,i)\\
    &+\Li_{6,1}(i,-1)+\Li_{5,2}(1,-i).
\end{align*}
For its integral shuffle, insert the single letter $-i$ into the word
$0^4(1)(-1)$ at each of its seven positions. Formula
\eqref{audit:s6-word-convention} then gives
\[
 P=\sum_{a=1}^{5}\Li_{a,6-a,1}(i,-i,-1)
       +\Li_{5,1,1}(1,i,i)+\Li_{5,1,1}(1,-1,-i).
\]
Subtract the two expansions and take imaginary parts. The first-index-one
terms have first argument $i$ and converge by summation by parts; all
remaining terms converge absolutely. Thus the finite-sum stuffle limit
and the convergent integral shuffle both apply. No divergent
$\Li_1(1)$ or regularization parameter is used.
\end{proof}

\paragraph{What this says about the known obstruction.}
The weight-seven functional $\ell$ in the pinned manuscript,
\texttt{cycloquot:eq:separating}, is a functional on a specified formal
depth-two presentation, not a functional on actual periods
\cite{Repo}. It annihilates the displayed single-by-single
shuffle/stuffle rows and the stated background; in particular,
\[
 \ell(K_{6,1})=1,\qquad
 \ell\bigl(\operatorname{Im}\Li_{5,2}(1,-i)\bigr)=0.
\]
Extend it to the free higher-depth word space by assigning zero to all
depth-three symbols, and call this particular extension $\widetilde\ell$.
For the convergent single-by-double row
\[
 R=\operatorname{Im}\{\mathsf{shuffle}(\Li_1(i),\Li_{5,1}(1,-1))
               -\mathsf{stuffle}(\Li_1(i),\Li_{5,1}(1,-1))\},
\]
the explicit expansions prove $\widetilde\ell(R)=-1$. Hence the naive
zero extension of the old separator already fails on this elementary
additional row. This calculation does not rule out another extension of
$\ell$, and it does not establish that any particular additional family,
including octahedral relations, is indispensable.

Let $Q_7$ denote the ten-term right side of
\eqref{audit:s6-depth-three-identity}. The ordinary stuffle of
$\Li_5(1)\Li_2(-i)$ gives
\[
 \operatorname{Im}\Li_{5,2}(1,-i)
 =g_{2,5}+\beta(7)-G\zeta(5),
 \qquad g_{a,b}=\operatorname{Im}\Li_{a,b}(i,1),\quad G=\beta(2).
\]
Selecting the even inner indices gives
$S_6=g_{6,1}+K_{6,1}$, so there is the exact reduction
\begin{equation}
 S_6=g_{6,1}-g_{2,5}+G\zeta(5)-\beta(7)+Q_7.
 \label{audit:s6-Q-reduction}
\end{equation}
The unresolved task is to reduce this specified $Q_7$ to the frozen
conjectural basket by a finite exact certificate. The double-shuffle
identities used here are classical; the contribution of this audit is
the explicit target row and the precise diagnostic of the inherited
obstruction.

The accompanying \texttt{verify\_s6\_target.py} independently enumerates
the seven integral shuffles and five stuffle terms using exact integer
arithmetic. It also checks all $192$ defining weight-seven restricted
rows against $\ell$, verifies the value $-1$ above, and records the
unchanged conjectural target in \texttt{s6\_target\_audit.json}.

% END sections/08a_s6_bridge.tex


\subsection{Three carry-forward manuscript corrections}

The following points were already identified in the incoming
Twisted Stieltjes report \cite{IncomingTwisted}. They remain relevant
at the pinned revision. The accompanying patch proposes only these
narrow changes and is supplied unapplied.

\paragraph{A real logarithm for the negative-parameter integral.}
In \texttt{chapters/03-algebraic.tex}, the displayed formula labelled
\texttt{cleo:eq:lintanh} uses $\log\tan(\theta/2)$ while its stated
real parameter range includes negative values. The real formula needs
$\log|\tan(\theta/2)|$. Under the principal complex logarithm,
the incorrect right side at $\lambda=-1/2$ differs from the real
integral by $-i\pi^2/6$. This is a branch correction, with the
absolute value fixed by the real integral and its real derivative.

\paragraph{Remove an unsupported nonreduction description.}
The same chapter calls a specific exact diagonal value
\emph{non-elementary} without defining a comparison algebra or proving
nonmembership. Keep the exact value and remove that classification.
An available special-function representation by itself proves no
nonreduction theorem.

\paragraph{Scope the coefficient field to the character in use.}
In \texttt{chapters/07-integration.tex}, the field
$\mathbb Q(\sqrt q)$ belongs to the particular quadratic-character
example. General additive or multiplicative character coordinates
can require cyclotomic and character values. The patch makes this
scope explicit without altering the quadratic calculation.

One earlier proposed wording change is already integrated: the PSLQ
discussion now allows known dependencies to be removed or quotiented
out and requires a nonzero coefficient of the target. It no longer
requires a prior linear-independence theorem. That edit is therefore
not repeated or counted as a new correction.

% BEGIN sections/08b_external_audit.tex
% The discussion concerns the inspected author PDF, not every published version.
% Bibliography keys to map in the parent: BorweinDilcher, DLMFpolylog.
\subsection{Corrections in the inspected Borwein--Dilcher author PDF}
\label{audit:BD}

The author-hosted PDF of Borwein and Dilcher,
\emph{Derivatives and fast evaluation of the Witten zeta function},
\url{https://carmamaths.org/resources/jon/omega3.pdf}, contains two
transcription errors relevant to a Mellin implementation.  Lemma~2.1(b),
printed page~2, assigns $H_0=1$; the required convention is $H_0=0$.
In Theorem~2.6, printed page~5, the positive-integer formula for $B_p$
has an extra factor $H_{p-1}$.  The corrected coefficients are
\begin{equation}\label{audit:BD-AB}
 A_p=\frac{(-1)^{p-1}H_{p-1}}{\Gamma(p)},\qquad
 B_p=\frac{(-1)^p}{\Gamma(p)},\qquad p=1,2,\ldots.
\end{equation}
These findings concern the inspected PDF; they do not assert that every
published version repeats them.

Here is an independent check.  Coalescing the Gamma pole and the
$\zeta(1)$ term in Jonqui\`ere's formula gives
\begin{equation}\label{audit:integer-Jonquiere}
 \operatorname{Li}_p(e^{-x})
 =\sum_{\substack{k\ge0\\k\ne p-1}}
       \zeta(p-k)\frac{(-x)^k}{k!}
 +\frac{(-x)^{p-1}}{(p-1)!}\bigl(H_{p-1}-\log x\bigr),
 \qquad 0<x<2\pi.
\end{equation}
Thus the coefficient multiplying $x^{p-1}\log x$ is precisely $B_p$
in \eqref{audit:BD-AB}.  At $p=3$, the required coefficient is $-1/2$,
whereas the printed $B_3$ is $-3/4$.  Separately,
$\operatorname{Li}_1(e^{-x})=-\log(1-e^{-x})=-\log x+O(x)$ forces
$H_0=0$.  A test at $p=2$ would miss the $B_p$ error because $H_1=1$.

Theorems~2.5 and~2.6 also state $\theta>0$ for their local power-series
expansions.  For ordinary convergent summation, impose
$0<\theta<2\pi$, as required by Jonqui\`ere's local series
\cite[25.12.12]{DLMFpolylog}.  A larger splitting point is possible only
with an additional continuation or intermediate-integral prescription.

To see the issue without an error estimate, take $r=s=0,t=3$ in the
double local series and keep $v=0,u=2k-1$.  The absolute value of this
summand is
\[
 \frac{|B_{2k}|\,\theta^{2k+2}}{2(2k)!(2k+2)}
 =\frac{\zeta(2k)\theta^2}{2k+2}
                       \left(\frac\theta{2\pi}\right)^{2k}.
\]
It fails to tend to zero for $\theta>2\pi$.  The original Witten sum at
$(0,0,3)$ converges, so the obstruction belongs to the displayed local
series representation.  The corrected implementation can retain the
classical theorem by choosing any $0<\theta<2\pi$; it does not need a
new assignment to divergent power series.

% END sections/08b_external_audit.tex


% BEGIN sections/08c_bailey_borwein_audit.tex
\subsection{Further transcription corrections in a Bailey--Borwein author PDF}
\label{audit:BB2014}

The following corrections concern the inspected author PDF of
\emph{Computation and theory of Mordell--Tornheim--Witten sums II},
dated 5 May 2014, printed pages 17--19 \cite{BB2014}. Equation numbers in
this subsection refer to that PDF. No assertion about other versions
is intended. Use its derivative convention
\begin{align}
 \omega_{a,b,c}(q,r,s)
 &=\sum_{n,m\ge1}
   \frac{(-\log n)^a(-\log m)^b(-\log(n+m))^c}
        {n^q m^r(n+m)^s},\notag\\
 \zeta_{a,b}(r,s)
 &=\sum_{k>j\ge1}\frac{(-\log k)^a(-\log j)^b}{k^rj^s}.
 \label{audit:BB-definitions}
\end{align}
Superscripts on $\Li$ and $\zeta$ below denote derivatives with respect
to their orders; $a,b,c$ are nonnegative integers.

\paragraph{Order-derivative signs in Eq.~(84).}
The factors printed as $(-1)^n\log^a n$ and $(-1)^m\log^b m$
must be $(-1)^a\log^a n$ and $(-1)^b\log^b m$.
Writing $\Gamma(s,z)=\int_z^\infty e^{-y}y^{s-1}\,dy$, the corrected
formula is
\begin{align}
 \omega_{a,b,c}(q,r,s)
 ={}&\sum_{n,m\ge1}\frac{(-\log n)^a(-\log m)^b}{n^q m^r}
     \partial_s^c\left\{
       \frac{\Gamma(s,n+m)}{\Gamma(s)(n+m)^s}\right\}\notag\\
 &+\int_{1/e}^{1}
    \partial_s^c\left\{\frac{(-\log u)^{s-1}}{\Gamma(s)}\right\}
    \Li_q^{(a)}(u)\Li_r^{(b)}(u)\,\frac{du}{u}.
 \label{audit:BB84}
\end{align}
Indeed $\partial_q^a n^{-q}=(-\log n)^a n^{-q}$.
Expanding the two polylogarithms in the integral from $0$ to $1/e$
and setting $y=-\log u$ gives the incomplete-Gamma term, including
the complete derivative of its Gamma ratio. There is no alternation
in $n$ or $m$.

\paragraph{The sign in Eq.~(86).}
The corrected identity is
\begin{equation}
 \omega_{0,0,b}(0,0,t)
 =\zeta^{(b)}(t-1)-\zeta^{(b)}(t).
 \label{audit:BB86}
\end{equation}
For $t>2$, exactly $k-1$ positive pairs have $n+m=k$, so its left
side is $\sum_{k\ge2}(k-1)(-\log k)^b/k^t$. This proves
\eqref{audit:BB86} directly. The printed right side has the opposite
sign; already at $b=0$ it is negative while the defining sum is positive.

\paragraph{Swapping orders also swaps derivative labels in Eqs.~(90)--(91).}
The corrected formulas are
\begin{align}
 \omega_{a,0,b}(s,0,t)+\omega_{b,0,a}(t,0,s)
 &=\zeta^{(a)}(s)\zeta^{(b)}(t)
        -\zeta^{(a+b)}(s+t),\label{audit:BB90}\\
 \omega_{a,b,0}(s,t,0)-\omega_{b,0,a}(t,0,s)
        -\omega_{a,0,b}(s,0,t)
 &=\zeta^{(a+b)}(s+t).\label{audit:BB91}
\end{align}
To see the indexing, $\omega_{a,0,b}(s,0,t)=\zeta_{b,a}(t,s)$.
Differentiate
$\zeta(t,s)+\zeta(s,t)=\zeta(s)\zeta(t)-\zeta(s+t)$
$a$ times in $s$ and $b$ times in $t$. The second omega term must
therefore carry the derivative labels $(b,0,a)$ when its spectral
arguments are $(t,0,s)$. The correction is invisible in the special
case $a=b$.

\paragraph{Gamma arguments in Eqs.~(93)--(94).}
Gamma functions and their derivatives are evaluated at the spectral
variable, not the number of differentiations. The corrected ladders are
\begin{align}
 \zeta_{b,a}(t,s)
 ={}&-\sum_{k=0}^{b-1}\binom bk
       \frac{\Gamma^{(b-k)}(t)}{\Gamma(t)}\zeta_{k,a}(t,s)\notag\\
 &+\frac1{\Gamma(t)}\int_0^1
   \frac{\Li_s^{(a)}(x)}{1-x}
   \log^b(-\log x)(-\log x)^{t-1}\,dx,
 \label{audit:BB93}\\
 \omega_{a,b,c}(q,r,s)
 ={}&-\sum_{k=0}^{c-1}\binom ck
       \frac{\Gamma^{(c-k)}(s)}{\Gamma(s)}\omega_{a,b,k}(q,r,s)\notag\\
 &+\frac1{\Gamma(s)}\int_0^1
   \Li_q^{(a)}(x)\Li_r^{(b)}(x)
   \log^c(-\log x)(-\log x)^{s-1}\,\frac{dx}{x}.
 \label{audit:BB94}
\end{align}
An empty sum is zero. For \eqref{audit:BB93}, differentiate
\[
 \Gamma(t)\zeta_{0,a}(t,s)
 =\int_0^1(-\log x)^{t-1}\frac{\Li_s^{(a)}(x)}{1-x}\,dx
\]
$b$ times with respect to $t$ and isolate the term with no derivative
on $\Gamma(t)$. The same Leibniz argument applied to
$\Gamma(s)\omega_{a,b,0}(q,r,s)$ gives \eqref{audit:BB94}.
Thus the occurrences of the Gamma argument $b$ in Eq.~(93) must be
$t$, and those of $c$ in Eq.~(94) must be $s$.

\paragraph{Direct convergence domains.}
These derivations require actual convergence before any continuation.
For real $q,r\ge0$, the integral formulas
\eqref{audit:BB84} and \eqref{audit:BB94} hold under
\begin{equation}
 s>0,\qquad q+s>1,\qquad r+s>1,\qquad q+r+s>2.
 \label{audit:BB-domain}
\end{equation}
The printed total-weight condition alone omits the two pair conditions:
at $(q,r,s)=(2,0,1)$ the $n=1$ subseries is harmonic and diverges.
Splitting the sum into $m\le n$ and $n<m$ proves sufficiency of the
three strict summability conditions; finite logarithmic powers do not
alter this open domain. Formula \eqref{audit:BB93} holds, for example,
when $t>1$ and $t+s>2$, while \eqref{audit:BB90}--\eqref{audit:BB91}
follow directly for $s,t>1$. The latter identities and
\eqref{audit:BB86} then extend meromorphically. Such continuation does
not assert ordinary convergence of the displayed integrals or series
outside their stated domains.

% END sections/08c_bailey_borwein_audit.tex


\subsection{Conventions that must survive integration}

No new false theorem was found among the inspected periodic Hurwitz,
coordinate-residue, and coincident-pair premises used in this article.
The following assumptions are nevertheless part of the mathematics
and must remain attached to the statements.

\begin{enumerate}
\item The pointwise product finite part, the periodic distributional
      convolution, and the canonical extension in the shift coordinate
      are different constructions. Theorem~\ref{cont:main} computes
      their relation, including its contact term.
\item The argument finite part uses a specified cutoff coordinate.
      In Theorem~\ref{cub:general-completion}, subtracting only the
      residue restricted to a polar hyperplane loses finite numerator
      derivatives and changes the requested answer.
\item A fixed-integer identity may be differentiated only in variables
      that remain free in that identity. In particular, the finite
      positive-integer Tornheim sums in \eqref{sp:positive} do not
      determine derivatives in either fixed inner order.
\item An origin ray is restricted before its removable singularity is
      filled in. The two distinct values $1/3$ and $5/12$ are both
      correct for their respective rays.
\item Numerical verification checks an implementation and can detect
      mistakes. It does not replace the proofs or establish arithmetic
      independence.
\end{enumerate}

The source audit and correction patch are separate from the new
mathematical section files. Their provenance notes identify which
findings are inherited, which are newly checked against an external
PDF, and which are restrictions on interpretation.

% END sections/08_audit.tex

% BEGIN sections/09_research.tex
\section{Further research questions}
\label{future:sec}

The formulas above suggest concrete next problems. The questions below
are proposals, not assertions of currently unproved identities.
They are ordered roughly by how directly the present proofs provide
the next calculation.

\subsection{Cubic and higher coincident products}

\paragraph{1. Reduce the third diagonal Tornheim derivative.}
Theorem~\ref{cub:diagonal} isolates $\omega'''(0)$ as the remaining
global coordinate for the cubic digamma finite part. Seek an exact
reduction of this derivative to a stated algebra of ordinary zeta
derivatives, Gamma or Barnes values, and convergent logarithmic
integrals. A useful first target is a second independent functional
identity for the cubic integral, rather than an unqualified
integer-relation search in a large constant basket.
If no such reduction is found, an explicit minimal integral
representation is still valuable, but it must not be called an
independence result.

\paragraph{2. Determine every asymmetric cubic ray derivative.}
In the local decomposition \eqref{sp:meromorphicH}, the next Taylor
layer involves $H_{AA}=H_{BB}$, $H_{AB}$, $H_{AC}=H_{BC}$, and
$H_{CC}$. The exact $C$ axis determines $H_{CC}$.
The diagonal cubic coefficient supplies one linear combination of
the other three. Can two additional exact restrictions or functional
equations determine the remainder? A concrete deliverable would be
a formula for $W_{a,b,c}'''(0)$ with its smallest explicit collection
of convergent-integral coordinates. This would clarify exactly where
ordinary single-zeta data cease to determine all directions.

\paragraph{3. Give the explicit global kernel at degree four.}
The local subtraction at any number of factors is already finite:
\eqref{cub:subset-numerator} lists the needed Taylor coefficients.
The unsolved part is the arithmetic of the global Fourier
zero-frequency constraint. At degree four, separate the sign
partitions of $n_1+n_2+n_3+n_4=0$ and identify a finite set of
constrained multiple sums whose meromorphic continuations represent
all moments. The first useful test is the fully coincident
$\FP\int_0^1\psi(x)^4\,dx$ with every counterterm written out.
This would extend the present cubic kernel without assuming
single-zeta closure.

\paragraph{4. Analyze triple collisions on unequal grids.}
The incoming pair formulas allow unequal spatial dilations.
A triple product introduces several collision loci and a possible
full three-way coincidence. Derive a chamber formula in the shift
variables, identify the exact changes on crossing each collision
wall, and compare the successive pair-collision limits with the
direct triple-coincident finite part. The challenge is to preserve
one stated coordinate prescription throughout; repeated finite-part
operations need not commute automatically.

\subsection{Contact distributions and coordinate transport}

\paragraph{5. Extend the contact generator to several shifts.}
Theorem~\ref{cont:main} determines the entire point-supported
correction for two factors and one shift. For a three-factor
correlation, the collision sets are lines and points in a two-shift
torus. Determine the distributions supported on those sets and their
intersections. The pairwise terms should recover
\eqref{cont:E}; the new information is the correction at a joint
intersection. A meaningful answer should specify test-function
normalizations and not introduce a product of singular
distributions without a definition.

\paragraph{6. Characterize invariant products exactly.}
Corollary~\ref{invg:stieltjes-products} gives a clean sufficient
criterion: at most one zero Stieltjes index. It is not a complete
classification of products or linear combinations of products.
Apply Theorem~\ref{invg:classification} to the finite Taylor jets of
the analytic factors and identify all cancellations among the
forbidden monomials. One concrete starting point is the span of
$\gamma_i\gamma_j\gamma_k$ at fixed total index, with coefficients
in a prescribed constant field. The field matters because a
functional cancellation and an arithmetic relation are different
questions.

\paragraph{7. Determine the coordinate-defect composition law.}
For two tangent coordinates $f$ and $g$, express the defect for
$f\circ g$ in terms of the defects for $f$ and $g$ and the
corresponding pullbacks. Identify a finite representation of this
law on pole order at most $R$ and logarithmic degree at most $L$.
The invariant subspace and its exact codimension are now known.
The remaining structural question is the action on the quotient,
including whether a natural coordinate choice gives a canonical
representative for each obstruction class.

\subsection{Harmonic Laurent data and primitives}

\paragraph{8. Move all the harmonic exponents independently.}
Theorem~\ref{h2:thm:direction} treats an outer exponent and the first
inner exponent while every trailing exponent remains one.
At the first new case
$\zeta_a(s,t_1,t_2)$, perturb all three orders independently.
Find the complete set of local polar divisors, then determine which
principal coefficients and finite parts admit finite Gamma,
Bernoulli, and Stirling formulas without an unevaluated
holomorphic remainder. General multiple-zeta Laurent algorithms
already exist \cite{h2:MOW}; the target is a compact explicit
specialization that preserves the common shift and all directions.

\paragraph{9. Classify directions with explicit higher regular terms.}
The two-exponent theorem stops at the finite part for a reason:
derivatives of its holomorphic remainder generally enter later.
The exact normal slice in Theorem~\ref{h2:thm:slice} is an exception
with an all-order Stieltjes formula. Determine other integer
outer-order slices or direction combinations for which the
remainder cancels to a prescribed order. A finite criterion would
organize which higher coefficients genuinely require new multiple
zeta derivatives.

\paragraph{10. Normalize primitives using Gamma multiplication.}
Equation~\eqref{h2:eq:primitive} supplies a finite primitive in the
shift variable, and \eqref{cont:Bernoulli-inverse} supplies
mean-zero periodic primitives. Study their compatibility with
reflection, multiplication of the shift, and rational shifts.
At low derivative orders, compare the resulting constants with
standard Barnes multiple-Gamma normalizations. Any claimed
simplification should specify both the additive constant and the
branch of the logarithm.

\subsection{Colored sums and exact Gaussian certificates}

\paragraph{11. Reduce mixed inner signs at arbitrary depth.}
The rational theorem settles the sector with all inner orders
nonpositive and arbitrary positive integer weights. Positive
integer orders at depth two are handled by partial fractions.
Develop a finite nested-polylogarithm reduction for arbitrary depth
with both positive and nonpositive inner orders, while retaining
the full complex outer order. Coincident colors and repeated poles
must be built into the formula. An affine shift alone is a routine
Lerch or Hurwitz extension and is less informative than this mixed
sector.

\paragraph{12. Detect arithmetic cancellation of depth-two jets.}
At roots of unity, \eqref{sp:DFTjets} and
\eqref{sp:DFTstieltjes} completely describe the depth-one terms.
Which character projections, symmetrizations, or distribution
combinations cancel the residual depth-two order derivatives in
\eqref{sp:11}? A useful result would be an exact finite linear map
whose kernel identifies combinations reducible to Gamma,
Hurwitz-zeta, and generalized-Stieltjes data. Such a map should
retain the allowed derivative directions instead of specializing
before differentiation.

\paragraph{13. Seek exact $S_6$ and $S_8$ certificates.}
Use the frozen target vectors at the pinned revision and enlarge the
proved relation space, not the numerical precision alone. The
existing quotient obstruction identifies relations that cannot
suffice; the verified $S_4$ octahedral proof identifies additional
families worth testing. For every added regularized relation,
prove the regulator prescription and retain its exact rational
coefficient row. A successful certificate must sum to the target
identity exactly. If a restricted search fails, publish the
precisely defined relation space and its separating functional,
without extending the impossibility statement beyond that space.

\paragraph{14. Formalize the finite proof kernels.}
The coefficient extractions in the subset completion, contact
generator, and harmonic residue polynomial are finite algebra.
Separate that algebra from the analytic lemmas of normal
convergence and meromorphic continuation, and formalize each part
with explicit interfaces. In particular, a reusable proof of the
Taylor-subtraction lemma in several spectral variables would
prevent accidental interchange of a coordinate finite part and a
spectral Laurent constant in future integrations.

% END sections/09_research.tex

% BEGIN sections/10_verification.tex
\section{Verification and reproducibility}
\label{verify:sec}

The proofs establish the all-parameter identities. The accompanying
code checks finite algebra, signs, specializations, and independent
analytic representations. Exact rational checks are distinguished
from arbitrary-precision floating-point diagnostics throughout.
The latter do not provide interval-certified error bounds.

\subsection{Independent algebraic checks}

The coordinate-invariance script compares two constructions of the
defect: a residue extraction and the binomial coefficient formula.
It performs 140 such exact comparisons, checks 36 obstruction ranks,
and verifies 440 Stieltjes threshold and sharp-coefficient instances.
All are finite symbolic computations.

The spectral script performs 402 exact checks: positive-inner-order
partial fractions, composition polynomials, harmonic--Stirling
coefficients, rational partial-fraction reconstructions, independently
convolved coefficient sequences, and the explicit weighted
cyclotomic certificate. The harmonic script verifies recovery of
the inherited one-direction residue formula for $m=0,\ldots,6$.
An independent symbolic expansion of the local Tornheim
decomposition gives identically zero residuals for the constant,
linear, and quadratic coefficients in arbitrary $a,b,c$.

For contacts, the direct $A$-generator and the symmetric
Gamma--cosine form \eqref{cont:B-symmetric} agree through total degree
four, after the exact even-zeta relations are applied. The script
also checks all seven displayed low-index coefficients and the
swap/reflection identity through total argument order five.

The focused $S_6$ audit independently regenerates the convergent
single-by-double row in Proposition~\ref{audit:s6-depth-three}, checks
all 192 rows of the restricted presentation against its separator,
and verifies the exact target normalization and two admissible
Cayley seeds. These checks certify the stated reduction and the
obstruction diagnostic; they do not certify the conjecture.

\subsection{Numerical checks using different representations}

\begin{longtable}{@{}>{\raggedright\arraybackslash}p{0.26\textwidth}
 >{\raggedright\arraybackslash}p{0.44\textwidth}
 >{\raggedright\arraybackslash}p{0.23\textwidth}@{}}
\toprule
Identity family & Independent comparison & Largest observed
absolute discrepancy\\
\midrule
\endhead
Cubic digamma finite part &
Local convergent digamma integral versus a
Jonqui\`ere/incomplete-Gamma evaluation of $\omega'''(0)$ &
$2.19\times10^{-35}$\\[3pt]
Multivariate cubic completion &
Direct Taylor-subtracted Hurwitz integral versus three independently
continued Tornheim kernels; real and complex asymmetric triples &
$5.84\times10^{-50}$\\[3pt]
Distributional contact generator &
Full meromorphic Hurwitz Fourier coefficient versus canonical
off-point coefficient plus contact; 189 mode comparisons &
$2.71\times10^{-76}$\\[3pt]
Quadratic Tornheim rays &
Independent Mellin continuation and symmetric extrapolation on
four distinct rays &
$2.60\times10^{-37}$\\[3pt]
Weighted and root-of-unity identities &
Nine convergent Mellin quadratures versus Gamma/zeta
specializations &
$4.68\times10^{-61}$\\[3pt]
Two-exponent harmonic Laurent data &
21 coefficients in eight cases, extracted by discrete Cauchy
averages from a continued nested sum &
$2.25\times10^{-46}$\\[3pt]
Stieltjes primitive hierarchy &
Six numerical shift derivatives compared with generalized
Stieltjes values &
$2.04\times10^{-56}$\\
\bottomrule
\end{longtable}

The cubic calculation uses 65 decimal working digits and refines
its finite-difference step. The leading known diagonal derivatives
are recovered well beyond the precision required for the cubic
comparison. The multivariate complex case was repeated with longer
independent truncations; its cross-error decreased to
$5.22\times10^{-58}$, while the direct integral changed by less
than $1.51\times10^{-63}$.

The contact check uses 85 working digits, three asymmetric complex
parameter pairs, frequencies $-3,1,4$, and all $p+q\le5$.
Its maximum scaled residual is less than
$2.52\times10^{-84}$. Omitting the proposed contact term in a
negative control produces an error of approximately $3.2457$,
so the comparison is sensitive to the delta contribution.
Zero modes and reflection signs are checked separately.

The harmonic extraction uses 65 working digits and keeps the first
64 original summands, with 32 Euler--Maclaurin tail terms and
32-point Cauchy averages on radius $0.01$. It does not subtract the
predicted Laurent expression before extraction. Repeating one
evaluation with cutoff 88 and 40 tail terms changes the result by
less than $5.9\times10^{-52}$. Details and parameters are preserved
in the JSON outputs.

\subsection{Package organization and integration}

The archive contains the modular \texttt{article.tex}, a flattened
standalone \TeX\ source, the compiled PDF, section sources,
reproducible Python scripts, their data outputs, and a source receipt.
It also includes a theorem-status ledger, independent review notes,
an integration guide, and the unapplied correction patch.
The \texttt{README.md} gives the build and verification commands.
The scripts depend on Python, SymPy, and mpmath; the document uses
standard pdf\LaTeX\ packages and requires no external bibliography
processor.

For manuscript integration, retain the proof of the cubic local
completion before the asymmetric ray theorem, because the ray proof
uses the diagonal quadratic coefficient. The contact and coordinate
classification sections can be integrated independently once their
finite-part conventions are retained. The spectral and harmonic
sections use independent notation and state the free variables
needed for their derivative claims.

The source receipt pins this continuation to a concrete repository
state. The patch was checked against that state and remains
unapplied. No theorem status in the repository is changed by the
existence of this archive; the result ledger is a proposal for
integration with the corresponding proofs.

% END sections/10_verification.tex


\begin{thebibliography}{99}
% BEGIN references.tex
\bibitem{Repo}
Vladimir Reshetnikov, \emph{ProveIt}, polylogarithm manuscript and
incoming research archives, revision
\texttt{dcd95baeb7c0e914b138bddef6829c5b1d52b726}.
\href{https://github.com/VladimirReshetnikov/ProveIt/tree/dcd95baeb7c0e914b138bddef6829c5b1d52b726/Analysis/Polylogarithms/docs/manuscript}{Pinned manuscript};
\href{https://github.com/VladimirReshetnikov/ProveIt/tree/dcd95baeb7c0e914b138bddef6829c5b1d52b726/docs/incoming}{pinned incoming directory}.
The accompanying \texttt{provenance/source\_receipt.json} records the
inspected source extraction.

\bibitem{IncomingCoincident}
ProveIt incoming report, \emph{Coincident Stieltjes Products},
10 October 2026, archive
\texttt{ProveIt\_Coincident\_Stieltjes\_2026-10-10.zip}.
\href{https://github.com/VladimirReshetnikov/ProveIt/blob/dcd95baeb7c0e914b138bddef6829c5b1d52b726/docs/incoming/ProveIt_Coincident_Stieltjes_2026-10-10.zip}{Pinned archive}.
Used for the quadratic completion, unequal-grid pair identities, and
the stated cubic and contact-term research questions.

\bibitem{IncomingTransport}
ProveIt incoming report, \emph{Exact Resonance and Coordinate
Transport}, 10 October 2026, archive
\texttt{ProveIt\_Resonance\_Coordinate\_Transport\_2026-10-10.zip}.
\href{https://github.com/VladimirReshetnikov/ProveIt/blob/dcd95baeb7c0e914b138bddef6829c5b1d52b726/docs/incoming/ProveIt_Resonance_Coordinate_Transport_2026-10-10.zip}{Pinned archive}.
Used for coordinate transport and the normal spectral-derivative
question.

\bibitem{IncomingExact}
ProveIt incoming report, \emph{Exact Identities}, 10 October 2026,
archive \texttt{ProveIt\_Exact\_Identities\_2026-10-10.zip}.
\href{https://github.com/VladimirReshetnikov/ProveIt/blob/dcd95baeb7c0e914b138bddef6829c5b1d52b726/docs/incoming/ProveIt_Exact_Identities_2026-10-10.zip}{Pinned archive}.
In particular, the fully shifted height-one theorem and the proposed
classification of coordinate-invariant logarithmic germs.

\bibitem{IncomingTwisted}
ProveIt incoming report, \emph{Twisted Stieltjes and Harmonic Laurent
Identities}, 10 October 2026, archive
\texttt{ProveIt\_Twisted\_Stieltjes\_Harmonic\_Laurent\_2026-10-10.zip}.
\href{https://github.com/VladimirReshetnikov/ProveIt/blob/dcd95baeb7c0e914b138bddef6829c5b1d52b726/docs/incoming/ProveIt_Twisted_Stieltjes_Harmonic_Laurent_2026-10-10.zip}{Pinned archive}.
Credited source of the carry-forward manuscript corrections.

\bibitem{EM2002}
Olivier Espinosa and Victor H.~Moll,
\emph{On some integrals involving the Hurwitz zeta function: Part 1},
The Ramanujan Journal \textbf{6} (2002), 159--188.
\url{https://www.math.tulane.edu/~vhm/papers_html/hurwitz1.pdf}.
Earlier preprint:
\url{https://arxiv.org/abs/math/0012078}.

\bibitem{EM2006}
Olivier Espinosa and Victor H.~Moll,
\emph{The evaluation of Tornheim double sums. Part 1},
Journal of Number Theory \textbf{116} (2006), 200--229.
\url{https://arxiv.org/abs/math/0505647}.

\bibitem{BB2014}
David H.~Bailey and Jonathan M.~Borwein,
\emph{Computation and theory of Mordell--Tornheim--Witten sums II},
author preprint dated 5 May 2014.
\url{https://www.davidhbailey.com/dhbpapers/mtw2.pdf}.
The audit in Section~\ref{audit:BB2014} concerns the equations and
convergence scope in printed pages 17--19 of this version.

\bibitem{BorweinDilcher}
Jonathan M.~Borwein and Karl Dilcher,
\emph{Derivatives and fast evaluation of the Tornheim zeta function},
The Ramanujan Journal \textbf{45} (2018), 413--432.
\href{https://doi.org/10.1007/s11139-017-9890-9}{doi:10.1007/s11139-017-9890-9}.
Inspected author preprint, titled \emph{Derivatives and fast evaluation
of the Witten zeta function}:
\url{https://carmamaths.org/resources/jon/omega3.pdf}.
The audit in Section~\ref{audit:BD} concerns this author PDF.

\bibitem{DLMFHzeta}
NIST Digital Library of Mathematical Functions,
\S25.11, \emph{Hurwitz zeta function}, especially the shift,
Bernoulli-value, Lerch, Fourier, and asymptotic formulas.
\url{https://dlmf.nist.gov/25.11}. Accessed 10 October 2026.

\bibitem{DLMFpolylog}
NIST Digital Library of Mathematical Functions,
\S25.12, \emph{Polylogarithms}, especially Eq.~25.12.12.
\url{https://dlmf.nist.gov/25.12}. Accessed 10 October 2026.

\bibitem{DLMFGamma}
NIST Digital Library of Mathematical Functions,
\S5.11, \emph{Asymptotic expansions}, especially Eq.~5.11.8
and the Gamma-ratio expansions in \S5.11(iii).
\url{https://dlmf.nist.gov/5.11}. Accessed 10 October 2026.

% END references.tex

% BEGIN references_contributions.tex
\bibitem{sp:Zhao}
Jianqiang Zhao, \emph{A note on colored Tornheim's double series},
Integers \textbf{10} (2010), A59, 879--882.
\href{https://doi.org/10.1515/integ.2010.059}{doi:10.1515/integ.2010.059}.
Author preprint:
\url{https://archive.mpim-bonn.mpg.de/1644/1/preprint_2009_66.pdf}.

\bibitem{sp:Romik}
Dan Romik, \emph{On the number of $n$-dimensional representations of
$\mathrm{SU}(3)$, the Bernoulli numbers, and the Witten zeta function},
Acta Arithmetica \textbf{180} (2017), 111--159.
\href{https://doi.org/10.4064/aa8455-3-2017}{doi:10.4064/aa8455-3-2017}.
Author version:
\url{https://danromik.com/resources/papers/su3reps-final.pdf}.

\bibitem{h2:AET}
S. Akiyama, S. Egami and Y. Tanigawa,
\emph{Analytic continuation of multiple zeta-functions and their values
at non-positive integers}, Acta Arithmetica 98 (2001), 107--116.
\href{https://math.tsukuba.ac.jp/~akiyama/papers/Mzffin.pdf}{Author-hosted text}.
DOI: \href{https://doi.org/10.4064/aa98-2-1}{10.4064/aa98-2-1}.

\bibitem{h2:MOW}
K. Matsumoto, T. Onozuka and I. Wakabayashi,
\emph{Laurent series expansions of multiple zeta-functions of
Euler--Zagier type at integer points}, arXiv:1601.05918.
\url{https://arxiv.org/html/1601.05918v1}.

\bibitem{h2:IshiiShinohara}
Y. Ishii and T. Shinohara,
\emph{Stirling polynomials and multiple zeta (star) functions at
non-positive integers}, arXiv:2508.12577; The Ramanujan Journal (2026),
DOI \href{https://doi.org/10.1007/s11139-026-01349-x}{10.1007/s11139-026-01349-x}.
\url{https://arxiv.org/abs/2508.12577}.
The comparison in this section is limited to the retrieved abstract.

\bibitem{Latt2011}
K. L\"att, \emph{The finite part of divergent integrals with logarithmic
factors}, Mathematical Modelling and Analysis \textbf{16} (2011), no.~4,
537--548. DOI: \href{https://doi.org/10.3846/13926292.2011.627951}{10.3846/13926292.2011.627951}.

% END references_contributions.tex

\end{thebibliography}
\end{document}
